\documentclass[../main.tex]{subfiles}
\begin{document}
%==========================================TIÊU ĐỀ TẬP========================================
\begin{table}[h]
    \begin{adjustwidth}{-1cm}{}
        \begin{tabular}{>{\centering\arraybackslash}p{6cm}|>{\raggedright\arraybackslash}p{12.5cm}}
            \multirow{5}{6cm}
            % Cột bên trái: ÉPISODE và số tập
            {%
                \\[-40pt]
                \flaregothic\fontsize{20pt}{20pt}\selectfont
                \centering
                \tikz\node{\contour{eptype}{\protect\color{black}\textbf{PERSONNALISÉ}}};
                \\[-8pt]
                \vnmdisp\fontsize{36pt}{36pt}\selectfont
                \tikz\node{\contour{epnum}{\protect\color{black}\textbf{HAI MƯƠI TƯ}}};
            }
             &
            % Dòng thứ nhất: tên tập tiếng Nhật
            {\fotskip\fontsize{18pt}{18pt}\selectfont ス◯◯チ２の\ruby{大}{だい}\ruby{虐}{ぎゃく}\ruby{殺}{さつ}}
            \\[6pt]
            % Dòng thứ hai: tên tập tiếng Anh
             & {\cafeteria\fontsize{18pt}{18pt}\selectfont The S*itch 2 Massacre}                                         \\[4pt]
            % Dòng thứ ba: tên tập tiếng Việt
             & {\vnmsans\fontsize{18pt}{18pt}\selectfont Xổ số S*itch 2}                                                  \\[8pt]
            % Dòng thứ tư: Nhân vật xuất hiện
             & {\sen{\fontsize{14pt}{14pt}\selectfont Track used: Battle Bastion (\ruby{戦}{せん}\ruby{闘}{とう}の\ruby{砦}{とりで})}
               }
            \\[30pt]
            % Dòng cuối cùng: Thời gian diễn ra của tập (thường sẽ là ngày phát hành)
            % & {\continuum\fontsize{16pt}{16pt}\selectfont Ngày phát hành: 07/02/2021}\\[18pt]
        \end{tabular}
    \end{adjustwidth}
\end{table}
%=========================================TRUYỆN CHÍNH========================================
\color{black}

\begin{center}
    \uline{Tháng Ba\\Văn phòng.}
\end{center}

Phòng thu âm của \hll\ hôm ấy đông vui bất ngờ. Dù không phải một chương trình phát sóng chính thức, nhưng chỉ cần nghe tin N*nt*nd* Direct sẽ công bố thế hệ máy chơi điện tử cầm tay tiếp theo, là bao nhiêu gương mặt đã háo hức kéo tới. Không khí bảo trùm như lễ hội.

Tất cả đều có mặt trước giờ G. Dàn ghế bành dài đặt trước màn hình lớn nhanh chóng kín chỗ. Pekora, Suisei, Ririka, Miko và Fubuki ngồi ở hàng đầum, chỗ được máy quay bắt rõ nhất. Ai cũng tay cầm đồ ăn vặt, miệng không ngớt bàn tán về chiếc máy cầm tay mới này.

Nàng Thỏ nhảy tưng lên trên ghế khi logo N*nt*ndo vừa xuất hiện: ``Cuối cùng thì ngày này cũng tới, peko! Switch 2 nhất định phải về tay tôi, trông nó hấp dẫn như thế nào luôn á!''

Nàng Ca sĩ vừa mở lon nước ra, vừa nghiêng đầu nhìn cô: ``Về tay em á? Với tốc độ bấm máy của em thì chị thấy là dễ về tay\dots\ người khác ấy. Nhưng mà công nhận nha, nếu lần này mượt thật thì M*r*o Kart chắc không còn cảnh tôi bay xuống vực vì lag nữa.''

``Cái gì chứ, peko!? Đừng có coi thường cô Thỏ này chứ! Kiểu gì may mắn cũng sẽ mỉm cười với tôi thôi, peko!''

``Ừ, may đến mức bị trượt khỏi đường đua tận ba lần liên tiếp kia kìa,'' Suisei đáp lại với tông giọng mỉa mai.
Cả nhóm bật cười, trong khi Pekora phồng má giận dỗi, không thể phản bác được gì. Ở đoạn phim giới thiệu đầu tiên, một phần trò chơi của Zelda mới hiện lên với đồ họa đậm chất điện ảnh, với từng chi tiết được thể hiện sắc nét.

Mọi người đồng loạt ``ooohhh'' lên, khuôn mặt ai nấy đều thể hiện rõ sự thích thú, thậm chí nàng Thầu khoán % Tên khác dùng cho Ririka, lấy từ "entrepreneur"
còn vỗ tay khá nhiệt tình, còn nàng Vu nữ thì cười đến tít mắt.

``Đồ họa đẹp vậy\dots\ chắc S*a*h Bros sẽ đã tay lắm! Cảnh va chạm mà đổ bóng tốt cỡ đó thì\dots\ trời ơi, chơi mê mẩn quên nấu cơm mất\dots'' đôi mắt của Ririka sáng lên, rồi bất giác nhìn tiền bối vu nữ: ``Mà mấy nhân vật nữ chắc còn sexy hơn nữa ấy chứ?''

Miko đáp lại với hai tay nắm chặt vào chiếc gối ôm hình bánh cá quen thuộc: ``Nyahaha\~{} Miko muốn chơi A*i** C******g trên cái máy đó thôi, ne! Mỗi lần tải đảo nhà là chị mất năm phút đó, chờ đến phát khóc luôn á!''

Trong khi đó, Fubuki - một game thủ chính hiệu - có hiểu biết sâu rộng hơn. Cô phân tích cặn kẽ với chiếc máy mới, dù không giấu nổi vẻ phấn khích: ``Chị còn nghe đồn rằng chiếc máy này sẽ sử dụng một loại DLSS đặc biệt dành riêng cho máy mới này đấy, lại còn có ray tracing luôn kìa!''

``Nếu mà Spl***n 4 lên được màn tốt nhất thì chắc là chúng ta sẵn sàng bán cả thận luôn ấy nhỉ,'' Suisei bông đùa. Các cô gái liền bật cười với câu nửa đùa nửa thật này.

``Mà này, Ririka! Nãy nói gì về mấy nhân vật nữ nhỉ? Bikini hay thân hình gợi cảm gì đó sao, peko?'' nàng Thỏ hất cằm nhìn về cô nàng hậu bối.

Ririka ngáp nhẹ, ngả người ra sau, khuôn mặt hiện rõ vẻ thỏa mãn: ``Ừ thì\dots\ nếu đẹp thì mặc ít cũng không sao mà? Đỡ nóng. Với lại em chơi game là để thưởng thức, không phải để giặt đồ ảo.''

``Vậy chắc chắn chị sẽ chơi trên con Switch mới này rồi, ne!'' Miko lanh lảnh. ``Chơi Spl***n kiểu gì hình đẹp, kiểu gì cũng nhìn rõ mắt hơn để trả thù cho Fubu-chan đó, ne!''

``Ừ nhỉ\~{}'' nàng Cáo trắng vểnh tai. ``Em còn giữ clip chị bắn trượt ba phát liên tiếp đó nha!''

``X-Xóa đi! Hãy đợi đó, chị sẽ phục thù em!'' nàng Vu nữ giơ nắm đấm lên cao tỏ vẻ thách thức với một trong bốn người chơi điện tử giỏi nhất, mặc dù khuôn mặt của cô có chút đỏ ửng vì xấu hổ.

Phản ứng có phần ngược đời này làm cho bốn cô gái còn lại tiếp tục có thêm một tràng cười nữa. Trong thâm tâm của họ, ai nấy đều mong muốn được sở hữu chiếc máy mới toanh này.

\begin{center}
    \uline{Nhà Hikawa\\8:00 tối.}
\end{center}

Một buổi tối sau khi cả gia đình Hikawa dùng bữa, ai nấy đều đã trở về  phòng của mình. Dưới ánh đèn huỳnh quang rọi xuống phòng khách ở khu A nhà Hikawa, một màn hình lớn đang chiếu lại đoạn ghi hình buổi N*t*nd* Direct từ chiều.

Ba thành viên của gia đình Tổng quản - Kuon, Kaigou và Miyako - mỗi người một vẻ đang xem lại buổi phát sóng trực tiếp trên tấm nệm được gập ba đặt ở góc tường.

Nàng Tổng ngồi khoanh chân, mắt dán vào màn hình với vẻ hứng thú. Cô vừa ăn bốc bánh tì tì ăn vừa bình luận như thể đang phân tích một bài toán chiến lược: ``\vie{Công nhận chiếc máy nhìn mượt thật\dots\ Bọn Nin bây giờ dồn lực vào trải nghiệm người chơi tốt hơn nha\~{}}''

``\vie{Cái đó anh phải công nhận đúng thật,}'' cậu Tổng ngồi ở bên kia tấm đệm đáp lại, nhưng ánh mắt cậu không nhìn về phía cô hay phía màn hình ti-vi. Thay vào đó, cậu đang đánh máy báo cáo đồ án bằng tốc độ của một người đã quen thuộc với chục năm nay.

``\vie{Em thấy là các hãng console kiểu gì cũng phải theo trend ray-tracing, upscale mấy cái model cũ cho người ta thấy là `tiến bộ'\dots\ Nhưng mà đẹp thật đó,}'' Kaigou tiếp tục, giọng có một chút hớn hở.

``\dots Ừm,'' nàng Chó trắng đáp nhẹ, mắt khẽ ngước lên màn hình ti-vi.

Nàng Tổng thở một tiếng khe khẽ, ngước đầu nhìn lên trần nhà: ``\vie{Ước gì mình có thể sở hữu một chiếc nhỉ.}''

Cô nàng vừa dứt câu, Kuon liếc thoáng qua màn hình vô tuyến rồi bĩu môi nhẹ: ``\vie{Anh thấy là cũng chẳng đáng lắm đâu.}''

``\vie{Sao lại thế?}''

``\vie{Thì có khác quái gì máy cũ đâu. Màn hình nét hơn, mạnh hơn tí, chứ cái lõi vẫn vậy. Có gì mà phải ầm ĩ ghê thế. Thậm chí cũng chỉ ngang một cái máy điện thoại tầm trung.}''

Giọng cậu hờ hững nhưng không giấu được sự thật rằng cậu còn hiểu biết hơn mọi thứ. Trong đầu Kuon, suy nghĩ còn rõ ràng hơn nhiều: ``\textit\vie{Thế nào rồi cũng có bản giả lập y chang trên PC. Nếu đã chơi single-player (chơi đơn) thì chạy trên laptop còn mượt hơn. Có mod, có save state, có upscale. Mua máy làm chi cho tốn.}''

``\vie{Aniki có vẻ như không hứng thú lắm với thứ này nhỉ?}'' cô em song sinh hỏi, cười trừ với phản ứng bất cần của cậu.

Cậu Tổng chỉ đáp lại bằng một cái lắc đầu, kèm với một từ ``Không'' đầy thẳng thừng, rồi tiếp tục tập trung vào màn hình laptop.

Kaigou liếc nhìn ông anh trai, không nói gì thêm, nhưng Miyako ngồi giữa hai người lại khẽ lên tiếng: ``\dots Ừm.''

Cô gái có đôi tai chó nhỏ và mái tóc trắng như sữa không nhìn TV, mà đang lật một cuốn sách minh họa các loài thực vật cổ. Nhưng rồi, cô ngẩng lên, mắt vẫn không chớp, nhẹ giọng: ``\vie{Máy mới thì chắc cũng không khác nhiều đâu\dots\ nhưng có máy thì vẫn dễ kết nối với mọi người hơn. Dù sao thì\dots\ nếu chơi cùng mọi người, thì máy nào cũng được.}''

Cô nói xong lại khẽ gật đầu, như thể tự xác nhận điều đó với bản thân, tiếp tục đắm chìm vào thế giới mộng mơ của mình.

Kuon dừng gõ một chút, tựa lưng vào tường, ngoảnh mặt về phía em gái: ``\vie{Nếu em thật sự muốn có một máy như thế thì sao không tự bỏ tiền ra mua đi?}''

Cậu ngừng một nhịp rồi tiếp lời, giọng điệu vẫn tỉnh rụi nhưng bắt đầu có phần cẩn trọng hơn: ``\vie{Anh nghe đồn là chiếc máy đó sẽ rơi vào khoảng năm man\footnote{Một man ở Nhật (1万円) là mười nghìn yên.} cho máy chỉ riêng cho Nhật, và thêm hai man nữa cho bản quốc tế, tức là đâu đó khoảng mười hai củ\footnote{Cách gọi lóng của người Việt. 1 củ là 1 triệu đồng.} Và em biết là nó bằng bao nhiêu không?}''

``\vie{Bao nhiêu?}''

``\vie{Hơn nửa kỳ học của anh em mình cộng lại đấy. Chưa kể nếu hàng hiếm thì nó còn tăng giá nữa cơ.}''

``\vie{Vậy em nghĩ chắc cái máy cũng đáng đấy chứ? Nó cũng chỉ cao hơn chút so với máy Một thôi mà,}'' Kaigou phản biện lại.

``\vie{Đấy là còn chưa nói đến giá game cơ. Hai củ cho mỗi trò, mà chưa chắc chơi được nhiều lần, thậm chí còn chẳng biết cái giá kia còn có đáng hay không cơ,}'' cậu nhún vai một cách nhẹ tênh, như thể con số đó chỉ là một đơn vị đo lường chứ không phải tiền thật.

Miyako vốn đang im lặng lắng nghe, lần này lại cất giọng chậm rãi: ``\dots ừm. Thế nên em thấy nó không đáng tiền là ở chỗ đó. Mặc dù tài chính của gia đình mình dư sức mua được, nhưng dùng số tiền đó để mua một món đồ cho trải nghiệm không quá khác biệt so với máy cũ thì em không thấy hợp lý.'' Ánh mắt lại lướt nhanh về phía Kaigou.

``\vie{Anh thấy máy Một cũng đủ tốt rồi mà,}'' cậu tiếp lời, hai tay cậu vẫn liên hồi gõ bàn phím. ``\vie{Mấy cái game mới của máy Hai có khi cũng chẳng khác mấy so với máy cũ lắm đâu.}''

Kaigou đang định phản biện, nhưng rồi lại lặng người một chút. Cô gác tay lên đầu gối, cầm điều khiển từ xa mà không bấm. ``\vie{\dots Ừm, cũng đúng. Mua một cái máy chơi game mà giá bằng cả một kỳ học thì\dots\ nghe nó cũng hơi bất thường thật.}''

``\vie{Tất nhiên là anh không cấm mua, tiền của em thì em muốn làm gì cũng được}'' Kuon lên tiếng, hướng về cô em gái đang có chút hơi hụt hẫng. ``\vie{Chỉ là anh nghĩ em nên suy nghĩ vấn đề này một cách thấu đáo hơn. Bây giờ cái gì cũng đắt, anh nghĩ ta nên tập trung vào đời sống thực tế hơn là cho game.}''

Nàng Chó trắng cũng chỉ gật đầu và ``ừm'' một tiếng, ngầm đồng tình với cậu.

Đáp lại họ, nàng Tổng chỉ thở dài một hơi. Cô lặng lẽ nhìn lại đoạn trailer Zelda mới trên màn hình, lần này không còn phấn khích như trước, mà có phần suy tính hơn.

Tưởng rằng ngọn lửa nhiệt huyết về chiếc máy mới đã tàn lụi sau những lời phân tích lạnh lùng của Kuon và kết luận thực tế từ Miyako, thì cuộc hội thoại ấy lại lọt vào tai hai người vừa chạy việc vặt của bố mẹ về: Ryuuhou và Suzuki.

Cửa vừa mở ra, Suzuki đã ló đầu vào phòng khách, giọng hồ hởi như thể vừa nghe tin có phát đồ ăn miễn phí: ``\vie{Ể!? Anh chị nói tới máy S*itch mới ạ!? Chơi M*r*o Party trên máy mới chắc vui lắm đó, onii-sama!}''

Ryuuhou chậm rãi tháo khẩu trang ra, đặt túi đồ lên bàn, rồi quay sang hỏi với ánh mắt có phần hứng thú: ``\vie{Máy mới hả? Có gì vui không, onii-san? Chắc là nó cũng tốt hơn đấy chứ?}''

Kuon ngả lưng, ngước lên trần nhà, đáp lại tỉnh bơ như đã lường trước: ``\vie{Ờ thì, mạnh hơn, màn đẹp hơn, chơi game mượt hơn. Nhưng mà giá chát. Chơi mấy cái đó giả lập trên laptop hay PC của anh cũng được rồi.}''

Cô em út cười toe, đáp trả lại: ``\vie{Nhưng laptop đâu có tháo rời tay cầm được đâu ạ! Với lại, S*itch chơi đông người nó vui hơn nhiều! Em chơi với mọi người là vui nhất á!}''

``\vie{Cái đó anh công nhận. Nhưng đó là việc của em, không phải của anh.}''

Ryuuhou theo đó cũng hùa theo: ``\vie{Nếu có game dễ thương nữa thì em cũng muốn chơi. Mình chơi chung với nhau cũng được mà?}''

Kaigou chép miệng nhẹ, vừa thấy tội cho hai cô em, vừa thấy họ nói có phần hợp tình hợp lý. ``\vie{Mua một máy cũng có thể dùng được cho cả nhà mà. Bọn em dùng là chính, không phiền gì anh đâu, aniki.}''

Dẫu cho những lời ngon ngọt của ba cô em gái, Kuon vẫn lắc đầu, giọng vẫn giữ sự điềm nhiên: ``\vie{Vấn đề đâu phải là ai dùng. Là có đáng không ấy. Đồ công nghệ mà nâng cấp nửa vời thì anh thấy phí tiền. Với lại-}''

Suzuki vẫn không hề bỏ cuộc, kéo tay áo cậu: ``\vie{Mua đi mà! Onii-sama là cool boy nhất nhà mà\~{}!}''

Cậu Tổng vẫn ngước mặt lên trần, làm lơ hoàn toàn lời cầu xin của cô em út, giơ hai tay lên cao như xin hàng: ``\vie{Việc đó em phải hỏi Kaigou ấy. Anh xin phép đứng ngoài vụ này.}''

Nhưng dẫu vậy, ba cô em gái nhà Hikawa vẫn tiếp tục gây áp lực lên cậu. Cô em song sinh của cậu ngồi bệt xuống sàn, nghiêng đầu nhìn cậu: ``\vie{Aniki à, em cũng thấy đáng mua lắm đó. Coi như là đầu tư giải trí gia đình.}''

``\vie{Thế thì cả nhà đầu tư luôn đi, anh không liên quan.}''

Trong lúc hai phía vẫn còn đang giằng co, Miyako vẫn ngồi xem sách tranh, nhấm nháp cốc trà sữa, rồi lên tiếng (dù không giúp ích lắm): ``\dots Em thì chơi máy gì cũng ổn. Có người chơi cùng là được.''

Chẳng mấy chốc, phe ``mua máy mới'' có ba thành viên: chính là ba cô em gái nhà Hikawa, đang cùng lúc tạo thành một tam giác tấn công quanh cậu cả. Nhưng dù bị dồn ép tứ phía, ông anh cả vẫn giữ vững lập trường, kiên định như tượng đá. ``\vie{Mấy đứa thích thì góp tiền vào mua nhé, anh không can tâm đâu!}''

Sau một hồi nói qua nói lại giữa hai phe nhưng không có kết quả, cuối cùng, cô chị song sinh của Kuon nảy ra một ý tưởng điên rồ nhưng rất hợp với không khí trong gia đình này. Cô vỗ tay một tiếng, làm bốn người còn lại chú ý: ``Vậy thế này. Mình chơi oẳn tù tì.''

``Ể?'' / ``Hả?'' Suzuki, Ryuuhou cùng Kuon nhíu mày.

``Nếu aniki thắng, bọn em sẽ không nói tới chuyện mua nữa. Nhưng nếu em, Ryuuhou hay Suzuki thắng\dots\ thì anh sẽ là người đi mua chiếc máy S*itch 2 kia. Không được từ chối.''

Cậu Tổng thở dài, chống cằm: ``\dots Oẳn tù tì chỉ để quyết một món đồ giá tận mười hai củ á?''

``Ừ. Bọn em biết dù sao anh cũng sẽ không mua nếu không có lý do hợp lý,'' cô nàng nói với giọng tinh nghịch: ``nên là bọn em chỉ còn có cách này thôi.''

Suzuki nắm tay cô em gái của cậu, phấn khích góp giọng: ``Chơi luôn! Em tin vào vận may của Ryuuhou-onee-sama luôn!''

Ryuuhou ngơ ngác, nhìn lại cô em út đang phấn khích kia: ``Ơ nhưng mà\dots\ chị còn chưa luyện chiêu\dots''

``Em sẽ tìm cách giúp chị đánh bại được anh ấy!''

Kuon thở dài, tặc lưỡi một tiếng. Biết rằng cậu không thể đọ được ba cô em gái đang hăng máu kia, cậu đặt chiếc laptop xuống, rồi nói: ``\dots Được rồi. Nhưng nếu anh thắng thì đừng bao giờ nhắc tới chuyện này nữa.''

``Hứa đi!'' cô em út chìa ngón út ra.

``Rồi, ngoắc tay. Anh không bao giờ hai lời đâu.''

\ct

Cả bốn người giãn ra thành vòng tròn nhỏ, những nắm tay giơ lên, chuẩn bị tinh thần ra trận. Không khí trong phòng như đông lại trong tích tắc.

Không ai cần nói thành lời, nhưng tất cả đều ngầm hiểu rằng: chỉ một keo duy nhất. Được ăn cả, ngã về không. Miyako, với chiếc cốc trà sữa còn lưng chừng trên bàn, tự động trở thành trọng tài, gật đầu như thể cô đã quá quen với việc này.

Kuon giữ gương mặt bình thản, đôi mắt nửa khép như thể vừa buồn ngủ vừa chẳng buồn để tâm. Nhưng trong đầu cậu, từng phân tích xác suất đơn giản cùng với những mẹo vặt cậu học trên mạng đang lướt qua như dòng mã chạy.

Kaigou đối diện với cậu đầu tiên. Hai ánh mắt giao nhau. Cô nghĩ rằng với sự hiểu nhau giữa hai anh em sinh đôi, mình có thể đánh lừa cậu chỉ một lần. Tay giơ lên, rồi hô:

``\dots Oẳn tù tì\dots\ ra cái gì\dots\ ra cái này!''

Kaigou ra kéo. Kuon ra búa.

Thất bại ngay ván duy nhất. Kaigou trề môi, ngồi phịch xuống tấm đệm. Cô nên biết rằng cậu sẽ đoán được kiểu ra tay ``trực giác'' của mình. Cậu ta đã sống cùng cô suốt ba năm, làm sao lại không biết cách cô sẽ suy luận ra sao chứ.

Tiếp theo là Ryuuhou.

Cô em gái hàng thật của cậu, đã sống với cậu từ thuở chập chững bước đi, còn lâu hơn cả Kaigou và Suzuki. Cô bé bước ra với vẻ mặt đầy quyết tâm, giơ nắm đấm của mình ra tư thế chuẩn bị.

Nhưng quyết tâm không thể thắng được những phân tích lạnh lùng của người anh trai.

``\dots Oẳn tù tì\dots\ ra cái gì\dots\ ra cái này!''

Giấy. Của cô. Kéo. Của anh trai cô.

Một lần nữa, chiến thắng nghiêng về phía cậu cả. Ryuuhou thở dài, mặt xị xuống nhưng vẫn nhẹ nhàng lui về sau, không trách móc gì. Miyako chỉ khẽ gật đầu xác nhận kết quả, không nói thêm một lời.

Tất cả hy vọng giờ đổ dồn vào Suzuki, cô em gái nhỏ tuổi nhất, nhưng lại là người liều lĩnh nhất và đôi lúc có những ý tưởng không ai ngờ tới.

Kuon khẽ liếc qua cô em, như đang tìm cách đọc ra biểu cảm từ ánh mắt, từ nhịp thở, từ cách đứng. Nhưng Suzuki lại chỉ đang\dots\ nghêu ngao hát thầm một giai điệu không ai hiểu, tay đung đưa như thể đây chỉ là một trò chơi trẻ con.

Như thể cô bé đang chọc tức anh trai mình vậy.

``\dots Oẳn tù tì\dots\ ra cái gì\dots\ ra cái này!''

Kết quả đã rõ. Nhưng lại theo chiều hướng ngược lại.

Kuon ra búa. Suzuki\dots\ ra giấy.

Im lặng bao trùm vài giây.

Suzuki đứng hình một chút, như không tin rằng mình vừa chiến thắng. Ngay sau đó, cả gương mặt cô sáng bừng lên như pháo hoa: ``\textbf{\vie{Em thắng rồi! Em thắng onii-sama rồi!}}''

Kaigou vỗ tay còn to hơn cả tiếng em gái mình hò lên điệu chiến thắng. Ryuuhou cũng reo lên khe khẽ, chung vui với cô em út của gia đình. Nàng Chó trắng chỉ ``ừm'' một tiếng, gật đầu xác nhận chiến thắng thuộc về cô bé nhỏ tuổi nhất.

Kuon thở dài, nhắm mắt ngửa đầu lên trần nhà, miệng mấp máy một cách mệt mỏi. Trong lòng cậu biết rõ: kể cả nếu hôm nay cậu thắng, ba cô em này cũng sẽ tìm cách khác để ép cậu ``thay mặt gia đình'' đi mua.

``\textit{Đằng nào thì cũng vậy cả. Mấy đứa này đúng là xảo quyệt thật.}''

Nhưng cái cảm giác bị đánh bại bởi đứa nhỏ nhất nhà trong trò oẳn tù tì vẫn khiến cậu ``cay đắng'' theo đúng nghĩa nhẹ nhàng của từ đó. Miyako sà tới bên cậu, vỗ vai như đang an ủi: ``\vie{Không sao đâu Cậu chủ. Đừng buồn.}''

``\vie{Onii-sama ơi\~{} Anh đã hứa rồi nha, mua máy cho bọn em đi nha\~{}}'' Suzuki buông ra câu hát đầy châm chọc, càng khiến cho Kuon thở dài mạnh hơn.

``\vie{\dots Thôi được. Quân tử nhất ngôn, đã hứa phải làm. Anh mua.}''

Chỉ cần nghe tới lời tuyên bố của cậu, cả ba cô em gái reo hò trong sung sướng, hí hửng khi cuối cùng họ cũng được cậu Tổng đồng ý mua chiếc máy.

\begin{center}
    \uline{Ngày hôm sau\\Văn phòng.}
\end{center}

Như thường lệ, văn phòng \hll\ đã rộn ràng tiếng cười nói từ sớm. Ánh nắng nhè nhẹ len qua rèm cửa, rọi xuống sàn nhà nơi những bước chân rôm rả đi qua. Kaigou bước vào, chiếc cặp chéo đung đưa theo nhịp đi, vẻ mặt không giấu được sự hào hứng. Đôi mắt cô ánh lên sự phấn khích như trẻ con vừa khám phá ra một món đồ chơi mới.

Theo sau là Kuon, cậu bước theo sau cô nàng đang nhảy chân sáo kia. Trái với sự hào hứng của cô em song sinh, gương mặt cậu hầu như không quá khác biệt so với thường ngày: vẫn điềm nhiên như không có gì xảy ra. Nhưng sâu trong đó, cậu vẫn còn hơi cay cú nhẹ vì cuối cùng cậu vẫn phải thực hiện mong muốn cho ba cô em gái của mình, dù thế cậu đã chấp nhận sự thật này rất nhanh và vì thế cũng không ảnh hưởng quá nhiều tới cậu.

``Chào buổi sáng mọi người\~{}!!'' giọng của nàng Tổng có chút lanh lảnh hơn thường lệ.

Pekora, Miko và Suisei đang ngồi ở bàn trên ghế bành, cùng đồng thanh đáp lại: ``Chào buổi sáng, Kaigou-senpai!'' Fubuki đang chỉnh lại bảng biểu trước khi buổi họp bắt đầu, còn Ririka thì vừa nhấm nháp cà phê vừa lướt lịch phát hành game, cũng ngước lên cười nhẹ.

``Có chuyện gì mà chị trông tươi tỉnh thế Kaigou-senpai? Hôm nay chị rạng rỡ thế,'' nàng Cáo trắng hỏi.

``Hehe\dots\~{} Mấy đứa biết gì chưa? N*nt*ndo Direct hôm qua công bố nhiều thứ hay lắm!''

Ngay khi cái tên N*n được nhắc đến, cả năm tài năng đều trông như bắt được vàng khi tìm được thêm một ``chiến hữu''.

``Có! Có chứ! Cái máy S*itch 2 đúng không!?'' nàng Vu nữ hớt hải hỏi.

``Chị đã xem cả cái chương trình hôm qua rồi chứ?'' nàng Ca sĩ tiếp lời, khuôn mặt không giấu nổi sự hứng thú.

Kaigou gật đầu, cười tít mắt: ``Và các chị biết sao không? Cả nhà chị sẽ đi rước S*itch 2 rồi đó!'' Như để chứng minh cho sự nhiệt tình về chiếc máy chơi điện tử mới, cô mở điện thoại vừa thao thao bất tuyệt về màn hình OLED, tốc độ tải nhanh hơn, thiết kế tay cầm mới, và còn phát lại các đoạn phim giới thiệu từ các trò chơi sẽ ra mắt trên máy mới.

Nàng Thỏ chen vào, đôi mắt sáng rực lên: ``Em cũng định mua chiếc máy ấy đó, peko! Hôm qua em coi lại ba lần cái trailer game mới! Nét gấp đôi máy cũ luôn nha, peko!!''

Nàng Vu nữ gật đầu liên tục. ``Miko cũng mê chiếc máy này ghê nhe\~{} Lúc chơi cùng mọi người mà còn mượt hơn nữa thì vui biết mấy!''

Ririka cười khẽ, tay vẫn đang cầm cốc cà phê: ``Có vẻ mùa hè này tụi mình sẽ có nhiều buổi stream multiplayer rồi đây.''

Trong khi tất cả đều ríu rít bàn tán, Kuon vẫn ngồi trước bàn làm việc quen thuộc. Màn hình laptop chiếu lên bảng cơ sở dữ liệu, đôi tai đeo một bên tai nghe, tay gõ những dòng lệnh ngắn để chuẩn bị chạy báo cáo hệ thống cho cấp trên.

Giống như những ngày khác. Im lặng, đều đặn, không bị cuốn theo những lời bàn tán của các cô gái về thứ đang nóng bỏng trong thời gian gần đây.

``Kuon-senpai\~{}'' Fubuki chợt gọi cậu từ phía bàn. ``Anh thấy sao về chiếc máy mới? Hôm qua anh có xem buổi ra mắt không đấy?''

Cậu ngẩng lên, tháo tai nghe ra. Một khoảnh khắc im lặng ngắn ngủi, rồi Kuon khẽ gật đầu: ``Ừ, cũng được. Công nghệ mới thì tốt thôi.''

Chỉ một câu gọn lỏn, không khen, không chê. Giọng điềm đạm, không gợn chút cảm xúc. Nhưng đủ lịch sự để không phá vỡ không khí vui vẻ trong phòng.

Kaigou liếc cậu một cái, nửa như buồn cười, nửa như trách móc ông anh bằng mắt: ``\textit{Anh thì lúc nào cũng vậy.}''

Nhưng các cô gái đều không hề biết rằng, trong thâm tâm cậu, cậu lại có một suy nghĩ khác có phần mạnh bạo hơn. Khi nghe từng lời tán dương, tâm trí cậu lại thoáng những so sánh không nói ra.

``\textit{Chip xử lý của máy Hai mạnh hơn\dots\ đúng, nhưng vẫn bị giới hạn vì phải tối ưu cho pin. RAM nhiều hơn thật, nhưng độ mở phần mềm vẫn kém, chưa kể rằng mình còn nghe đồn chiếc máy mới còn hạn chế khả năng mở rộng nữa. Giá đó dựng được cái PC mạnh gấp ba lần. Mà giả lập thì có thiếu game gì đâu?}''

Cậu nhớ lại dàn máy cây ở nhà, thứ mà cậu từng lắp cùng Koheki - một người bạn hồi cấp ba giờ đang học ở một trường quân sự, với CPU nhà Đỏ\footnote{Ám chỉ đến AMD.} và GPU tầm trung nhưng tối ưu cực tốt, có thể chạy giả lập cho máy S*itch với tốc độ chẳng thua gì máy thật.

Nhưng Kuon không nói ra. Không phải vì ngại, mà vì cậu thấy không cần thiết. Ai cũng có niềm vui riêng, và đôi khi, để yên cho nó tồn tại cũng là một cách để giữ hòa khí. Nhất là lần  cậu đã vô tình ``nói xấu'' về món bánh pudding, để rồi bị Marine, Haato và Choco mắng xối xả và ``tặng'' cho cậu vài bài thuyết giảng về sự tuyệt vời của nó như thể họ đang ở một nhóm tà đạo\footnote{Xem lại {\flaregothic ÉPISODE 50}, quyển số Bốn.}.

Thế nên cậu chơi bài dù an toàn mà được lòng mọi người hơn: mỉm cười nhẹ, tay lại đặt lên bàn phím, tiếp tục công việc như mọi khi, một cách lặng lẽ và ổn định.

Câu chuyện quanh bàn trà giữa các cô gái vẫn rôm rả, nhưng ánh mắt của nàng Cáo trắng thi thoảng lại lướt về phía Kuon, người duy nhất không cùng họ nói về chiếc máy mới. Ánh sáng buổi sáng chiếu nhẹ lên đôi mắt đỏ của cậu, tạo nên một lớp phản chiếu nhẹ khiến biểu cảm càng khó đoán hơn thường lệ.

Fubuki nghiêng đầu, chống cằm nhìn Kaigou, cô em gái vẫn đang thao thao bất tuyệt về ``dock mới có hỗ trợ 4K, J*y-C*n đỡ bị trượt hơn, và nghe nói có game Spl*t**n kiểu battle royale (đấu trường).''

``Hmm\dots\ Kaigou-senpai này,'' nàng Cáo trắng chậm rãi mở lời, giọng có chút tinh nghịch nhưng cũng đầy ẩn ý. ``Chị thì háo hức như vậy, mà sao Kuon-senpai nhìn như không quan tâm luôn. Hai người là anh em ruột thật không đó?''

Cả bàn cùng phá lên cười. Nàng Tổng đang uống nước, suýt thì sặc. ``Thật mà! Nhưng tính anh ấy vốn thế!'' cô phân bua, mặt hơi đỏ lên.

Miko liếc qua Kuon, rồi cười toe: ``Em thử đoán nha\~{} Chắc là anh ấy ngại thể hiện thôi ấy mà, ne. Đúng hơm nào, Kaigou-senpai?''

Suisei nhướng mày, quay sang Ririka đang im lặng quan sát nãy giờ: ``Ririka-chan nghĩ sao? Kuon-senpai đang cố giấu phấn khích chăng?''

Ririka nhún vai, cười khẽ: ``Khó nói lắm. Kiểu người đó nếu có hứng cũng không lộ ra mặt đâu. Phải ép ra thì may ra mới khai.''

Pekora đập tay xuống bàn một cái, giọng lớn rõ: ``Vậy thì hỏi thẳng luôn, peko!'' Không nói gì nhiều, cô nàng Thỏ lon ton tiến tới cậu con trai vẫn đang lập trình. ``Kuon-senpai, anh biết gì về cái máy mới không? Nói thật đi!''

``Cái gì nữa?'' cậu hỏi vọng lại. ``Lại dạng giống câu hỏi thử lòng đấy hả?''

Theo sau cô, bốn tài năng còn lại cũng tiến tới, đồng thanh: ``Đúng đó! Nói thật điii\~{}!''

Bị bao quanh bởi một vòng vây đầy ``nghi vấn'' từ năm cô gái, Kuon ngẩng lên khỏi màn hình, gỡ tai nghe xuống lần nữa, bình thản như thể đang trả lời một câu hỏi kỹ thuật.

``Biết chứ. Kiến trúc vi xử lý mới, chip mới nhà NV*DIA tăng cơ chế bảo mật, hệ thống dock có thể upscale lên 4K, có nút C cho phép giao tiếp tốt hơn cùng với phụ kiện camera. Có những tựa game sắp tới gồm M*i*o Kart World hay The Legends of Z*l*a mới, F*nal F*nt*sy bản làm lại, vân vân, rồi có tin đồn Fire Emblem mới. Máy có thể lên tới 12GB, bộ nhớ trong 256GBm màn hình 7.9 inch, màn LED. Có hai phiên bản khác nhau: một bản nội địa Nhật giá khoảng 5 vạn yên và bản quốc tế giá 7 vạn yên, chưa tính thuế.''

Cả năm cô gái đều ``ồ'' lên. ``Quả không hổ danh Kuon-senpai có khác, đúng là anh có xem buổi ra mắt đó, ne!''

Cậu thở dài, rồi tiếp tục vẫn với giọng điềm đạm: ``Nhưng mà biết mấy thứ đó không có nghĩa là anh thích nó, Miko-chan ạ.''

Năm cô gái nhìn nhau trong thoáng chốc, rồi cùng gật gù như vỡ lẽ. Nàng Thầu khoán chống nạnh: ``Thế thì anh không phải là kiểu `giả vờ ngầu lòi trong lòng thì phấn khích' à?''

``Không,'' cậu đáp lại ngắn gọn, không chút chần chừ.

Nàng Cáo trắng bật cười lớn: ``Ái cha cha\~{} Thế mà em tưởng anh là kiểu tsundere công nghệ chứ!''

Kaigou bật cười ngặt nghẽo, lấy tay đập nhẹ vào bàn: ``Tsundere công nghệ là cái gì vậy trời!? Em tự dưng sáng tác ra thêm định nghĩa mới hả, Fubu!''

Trong lúc mọi người đang cười đùa, Kuon chỉ nhún vai, đeo tai nghe lại. Cậu không giận, cũng không xấu hổ, cậu chỉ đơn giản là không thấy lý do phải thể hiện điều gì nếu không thực sự cảm thấy thế.

Với cậu, cái máy kia vẫn chỉ là một món đồ đắt đỏ. Và khi nhìn từ góc độ kỹ thuật, nó vẫn thực sự chưa ``đáng tiền'' như mọi người tin tưởng. Nhưng nếu em gái mình và những người cậu quý trọng vui vẻ vì điều đó, thì cậu sẵn sàng để bản thân đứng ngoài niềm vui ấy, chứ không chống đối lại nó.

Ngay khi tiếng cười nói vừa dịu lại, Fubuki đột ngột chuyển tông, vỗ hai tay vào nhau một cái *BỐP!* rõ to, khiến mọi người giật mình. ``À mà nè! Nói vậy chứ muốn mua S*itch 2 bây giờ không phải dễ đâu nha!''

Miko nghiêng đầu: ``Ể? Ý em là sao, ne?''

``Xổ số đó!'' nàng Cáo trắng nhấn mạnh, giọng đầy kịch tính. ``Phải đăng ký xổ số mới được quyền mua nha!''

Một nhịp lặng kéo dài đúng ba giây.

``Nghĩa là\dots'' Suisei mở to mắt, ``\dots không phải cứ có tiền là mua được à!?''

``Ừm.'' Fubuki nghiêm túc gật đầu. ``Đợt này chỉ có năm trăm nghìn máy bán ra, mà lượt đăng ký lên tới hai triệu người rồi. Tỷ lệ khoảng 25\%. Mà chưa hết đâu, đầu giờ chiều nay là hết hạn đăng ký rồi đó!''

Cả Miko, Suisei, Pekora và Ririka đồng loạt đổ mồ hôi hột. Họ không nói gì, chỉ nhìn nhau bằng ánh mắt đầy rùng rợn, lẩm bẩm: ``Nãy giờ mình cứ nghĩ là dễ mua chứ\dots''

Không ai bảo ai, họ cùng rút điện thoại ra đăng ký ngay, tốc độ còn nhanh hơn lúc bị rượt vì chậm tiến độ. Nàng Thỏ vừa thao tác vừa lẩm bẩm: ``My N*nt*ndo ID\dots\ mật khẩu\dots\ Không nhớ, peko- khoan đã, lưu đâu ta!?''

``Bà chơi từng ấy năm mà còn không nhớ nổi mật khẩu ư, ne!?'' Miko thảng thốt, rồi giở giọng châm biếm: ``Vậy có khi bà không phải là fan của N*nt*ndo chính hiệu rồi\~{}''

``I-Im đi!''

Ririka thì nhanh chóng đăng nhập, tay lướt điêu luyện: ``Em đăng ký rồi nhé. Cho chắc.''

Kaigou lúc đầu còn ngơ ngác, nay thì bối rối hẳn. Cô đảo mắt qua lại giữa mọi người, rồi hỏi như không tin vào tai mình: ``Khoan đã khoan đã\dots\ Là thật hả? Chơi xổ số thật á!? Không phải đùa hả Fubu!?''

Suisei ngẩng lên, lần này giọng nghiêm túc hơn thường lệ: ``Không đùa đâu Kaigou-senpai. Nếu chị không đăng ký trong sáng nay, cơ hội lần này coi như đi luôn đó.''

Miko cũng không còn nụ cười thường ngày, sắc mặt dần chuyển sang sự căng thẳng tột độ: ``Lần này người ta giới hạn đăng ký rất gắt luôn ne\dots\ Phải chơi ít nhất 50 tiếng trên Switch, và phải có tối thiểu một năm đăng ký thành viên cho NSO nữa ne!''

``Và tài khoản NSO của chị cũng phải đang hoạt động nữa đó!'' Ririka chêm thêm.

Kaigou cắn môi, lật đật rút điện thoại, tay run run. Không phải vì hoảng loạn, mà là vì không ngờ chuyện ``mua máy chơi game'' lại đột ngột trở nên rối rắm như thế này. Cô vừa nhập mật khẩu vừa hỏi: ``Mà nếu mình không trúng thì sao\dots?''

Nàng Ca sĩ nhún vai: ``Thì đợi đợt sau.''

``Mà đợt sau có khi phải qua tháng sau luôn đó!'' nàng Vu nữ chêm thêm, không quên cúi đầu sát tai thì thầm như kể chuyện ma.

Và thế là, họ không ngờ rằng, việc rinh về một chiếc máy chơi điện tử mới lại không hề đơn giản như vậy. Thậm chí có thể nói là khốc liệt không kém gì giành được một vé vào cánh cửa đại học top đầu.

Trong lúc mọi người rối rít, bầu không khí trong căn phòng như đặc quánh lại vì áp lực vô hình từ chiếc máy bé nhỏ ấy\dots

\dots chỉ có một người vẫn ngồi yên bất động, cũng là người không tham gia cuộc trò chuyện ban đầu ấy. Cậu không nói gì, không bình luận, chỉ thoáng nhếch mép, nụ cười mỉa mai như thoáng qua gương mặt. ``\textit{Đến mua một cái máy chơi game cũng phải phụ thuộc may rủi. Thật là hết thuốc chữa.}''

Cậu không trách ai, cũng không trách thời cuộc. Chỉ đơn giản là thấy cái xã hội này đôi lúc thật lạ kỳ. Nhưng nếu đó là cách để mọi người cảm thấy ``có hy vọng'' hơn, thì cậu cũng chẳng định phá vỡ điều đó.

Giữa lúc văn phòng vẫn còn sôi nổi vì chuyện đăng ký, Kuon mới nghiêng người về phía Kaigou, khẽ nói, giọng như chỉ muốn cô nghe: ``\vie{Nếu không mua được lần này thì lần sau mua cũng được. Em không cần quá lo đâu.}''

Cô em song sinh của cậu liếc nhìn bằng ánh mắt kiên định: ``\vie{Lần này là được thì phải cố, aniki ạ. Đã bao lâu rồi em mới có chút hứng thú thật sự với thứ gì đó ngoài giấy tờ và mail đó.}''

Cậu Tổng thở dài, tiếp tục khuyên nhủ: ``\vie{Anh hiểu ý của em rồi, nhưng mà nhiều khi ấy, hàng lượt đầu chưa chắc đã có chất lượng tốt đâu. Còn phải đi khảo sát thị trường rồi còn nghe ngóng người ta đánh giá nữa. Lúc đó quyết định mua hay không còn chưa muộn.}''

``\vie{Thế anh không nghĩ tới cảnh Suzu sẽ buồn vì không giữ lời hứa sao?}'' Kaigou hướng khuôn mặt nghiêm túc về ông anh trai. ``\vie{Anh còn nhớ chuyện gì xảy ra với con bé tháng trước khi ta thất hứa chứ?}''

Cậu nhớ lại cảnh hồi vài tháng trước, khi Kuon và Kaigou có hứa sẽ mua kem cho Suzuki khi họ đi làm về, nhưng rồi bởi vì họ quên, thế là cô em út la lối om sòm lên, chạy rượt khắp cả khu nhà, làm náo loạn cả bốn khu, tới mức\dots\ suýt làm căn nhà nổ tung vì quá tải.

``\vie{\dots Cũng phải. Nó cũng cứng đầu lắm chứ,}'' cậu cười nhạt, tỏ vẻ không muốn nhớ đến ký ức kinh hoàng đó. ``\vie{Hồi đó không có Ryuuhou với Miyako can ngăn thì có khi nhà ta còn mỗi cái nịt thôi. Nó cũng gần hai mươi tuổi đầu rồi còn gì.}''

``\vie{Thông cảm cho Suzu đi, aniki ạ\dots\ Lâu lắm rồi nó mới được sung sướng như vậy đấy. Ít nhất cũng để cho nó tận hưởng nó chứ.}''

``\vie{Em mà chiều nó thế cẩn thận hư người đấy.}''

``\vie{Biết rồi, nói mãi.}''

Hai anh em đứng tách khỏi hội năm cô gái, vừa nói chuyện vừa lướt qua các thông tin chi tiết về đợt xổ số. Không lâu sau, họ bắt đầu bàn tính chiến lược, như thể đang lên kế hoạch cho một nhiệm vụ mật.

``\vie{Theo anh, ta sẽ đăng ký bằng hai tài khoản,}'' cậu Tổng đề xuất. ``\vie{Một của anh, một của em. Dù sao thì cơ hội càng cao càng tốt.}''

``\vie{Tài khoản của anh là từ cái máy Sui-chan tặng hôm sinh nhật hồi ba năm trước phải không?}'' cô em hỏi lại.

``\vie{Ừ. Cái máy đó anh dùng được khoảng ba tháng, rồi bàn giao cho Ryuuhou dùng tiếp. Đến giờ vẫn dùng bình thường. Còn em thì\dots}''

``\vie{Máy em là do trúng thưởng bốc thăm ở Comiket hồi 2023. Chưa đổi.}''

Cả hai cùng nhìn nhau gật đầu, yêu cầu tối thiểu của họ là phải có một chiếc máy coi như đã xong. Tiếp theo, cả hai người họ cùng nhau đọc các điều kiện để kiểm tra xem điều kiện được tham dự xổ số là gì.

``Xem nào\dots'' cô em song sinh nhíu mày, đọc những dòng thông tin trên trang chính thức của nhà Nin. ``\vie{Điều kiện đầu tiên là phải \emph{tối thiểu 50 giờ chơi.}}''

Kuon lẩm bẩm, nhẩm nghĩ: ``\vie{Anh thấy Ryuuhou ngày nào cũng chơi ít nhất hai ba tiếng. Nó còn mang máy đi học cả ngày thứ Bảy tuần trước nữa. Chắc chắn thừa 50 giờ.}''

Kaigou cũng xác nhận: ``\vie{Máy em thì lâu lâu mới mở, nhưng đợt này để trong phòng khách nên Mi-chan với Suzu chơi suốt. Chắc kèo là không dưới trăm giờ.}''

``\vie{Như vậy là anh em mình đạt được điều kiện thứ nhất rồi. Để xem cái tiếp theo là gì\dots}'' cậu lướt xuống thêm một hồi, dừng lại ở một dòng khác: ``\emph{Tài khoản của người dùng phải đăng ký gói hội viên NSO tối thiểu 1 năm.}''

Cậu Tổng nhướng mày: ``\vie{Anh không rõ em đăng ký membership thế nào\dots\ Em còn nhớ không?}''

Nàng Tổng đáp ngay: ``\vie{Em mua luôn gói năm từ lúc cài máy, nên chắc chắn là đến giờ được hơn một năm rồi.}''

Kuon khoanh tay, thở ra một tiếng. ``\vie{Lúc chuyển máy cho Ryuuhou vẫn giữ NSO trên tài khoản chính, anh còn bảo nó là cứ chơi trên nick của anh đi. Giờ vẫn còn hiệu lực.}''

Cả hai anh em liếc xuông tới dòng thứ ba, cũng là điều kiện cuối để tham dự: ``NSO vẫn còn hiệu lực vào thời điểm đăng ký.''

Cậu con trai nhớ lại, gõ nhẹ vào trán: ``\vie{Anh nhớ là cuối tháng Tư này mới phải làm mới lại. Vẫn dư hơn một tháng.}''

Cô em song sinh gật đầu, hừm một tiếng như đạt được mong muốn của mình. ``\vie{Của em là tháng Năm. An toàn.}''

Như vậy, anh em nhà Hikawa coi như đã đạt được cả ba điều kiện được coi là cần để có một suất tham dự xổ số máy Hai. ``\vie{Vậy là ổn. Mỗi đứa một suất, tăng gấp đôi cơ hội.}''

Đến cuối cùng, chuyện tưởng chừng như khá nhỏ nhặt, lại có thể rất kỳ công như thể đang lập kế hoạch cho bài đồ án tốt nghiệp. Sau khi chắc chắn về điều kiện, họ bắt đầu chọn phiên bản máy.

Kaigou nhìn hai ô chọn trên màn hình, lia qua lia lại con chuột giữa chúng: ``\vie{Nội địa hay quốc tế đây\dots}''

Kuon đứng cạnh cô, khoanh tay suy nghĩ một chất tư vấn: ``\vie{Lấy cho anh bản quốc tế. Đắt hơn hai vạn yên, nhưng anh thấy người ta thường có xu hướng chọn bản nội địa do rẻ tiền và hợp với như cầu hơn.}''

``\vie{Anh nghĩ như vậy á?}''

``\vie{Suy luận xã hội một chút là ra mà. Anh cũng thấy mấy đứa quen đọc tiếng Anh hơn là tiếng Nhật, nên chắc là sẽ hợp gu hơn. Vậy nên, nhắm vào bản càng ít người mua càng có cơ hội chiến thắng.}''

``\vie{\dots Cũng phải.}''

Sau khi chọn xong, hệ thống hiển thị Mã dự thưởng. Kuon nhìn chằm chằm vào bốn số cuối trong dãy mã của mình: 4953.

``\vie{Hmmm\dots\ Bốn chín năm ba. Vãi chưởng.}'' Cậu nhăn mặt.

Kaigou nhìn số đó, thoáng run nhẹ: ``\vie{Đã 49 còn 53 nữa. Xui gấp đôi luôn kia, aniki ơi\dots}''

Cậu Tổng cười nhạt với dãy số đen đủi kia, tự trấn an bản thân: ``\vie{Cũng chưa chắc đâu. Anh nghe nói có một vài người bấm trúng biển số xe dạng này vẫn gặp may mắn đấy thôi. Dù gì cũng chỉ là số thôi mà.}''

Đến lượt Kaigou, màn hình hiện lên dãy mã có bốn số cuối là 8888. ``A\dots''

``\vie{Tứ bát cơ à?}'' cậu nhướng mày. ``\vie{Em số đỏ phết.}''

Kaigou đổ mồ hôi mà vẫn không nhịn được cười: ``\vie{Kiểu này dễ ăn may lắm đó. Làm ơn cho em đậu lần này đi\dots}''

Cả hai cùng im lặng một lúc. Chỉ còn tiếng bấm phím rì rào và không khí kỳ lạ giữa hai con số đối lập đến khó tin: một bên là điềm rủi trùng trùng, một bên là vận đỏ bốn phương.

\ct

Mười một giờ trưa.

Tiếng chuông đồng hồ trên điện thoại Fubuki vang lên đúng lúc mọi người vừa hoàn tất các bước đăng ký cuối cùng.

``\textbf{Hết hạn rồi đó, mọi người!}'' nàng Cáo trắng hô lớn, tay giơ cao điện thoại, đồng thời ngã người ra ghế như vừa hoàn thành một nhiệm vụ sinh tử. ``Ai chưa bấm nút gửi là tiêu đời rồi nhaaa\~{}!''

Cả văn phòng như vừa thoát khỏi một cơn bão. Nàng Tổng thì vẫn còn đang nhìn màn hình, thở phào nhẹ nhõm khi thấy thông báo ``Đã hoàn tất đăng ký'' hiện lên. Miko thì ôm ngực thở hổn hển, còn Suisei thì ngồi phịch xuống ghế, vung tay: ``Vậy là xong! Một vé cược vận mệnh!''

Ririka huơ huơ tay: ``Em có số đẹp nè! Số em là\dots\ 1020! Có ai giống em không?''

Pekora nhảy lò cò quanh chỗ ngồi, tai thỏ đung đưa: ``Tôi mà trúng thì sẽ livestream unbox liền ngay khi nhận hàng nha, peko!''

Fubuki đảo mắt một vòng, tay chống hông: ``Nào, nào! Báo cáo tình hình: Miko-chan?''

``Xong luôn rồi ne!''

``Sui-chan thì sao?''

``Rồi! Số em đẹp luôn!''

``Ririka?''

``Xong rồi ạ!''

``Kaigou-senpai?''

``Chị cũng xong rồi đây\dots\ Hú hồn!''

Nàng Cáo trắng gật đầu mãn nguyện. Mọi ánh mắt liếc về phía Kuon, người từ đầu tới giờ vẫn chẳng thèm quay đầu lại. Tay vẫn gõ bàn phím đều đặn, mắt dán vào màn hình với vẻ dửng dưng quen thuộc.

``\dots Kuon-senpai?'' cô nheo mắt, hỏi với tông giọng nghịch ngợm.

Cậu vẫn không quay đầu, chỉ đáp gọn: ``Anh chỉ giúp Kaigou thôi. Xong rồi thì làm việc tiếp. Còn đầy việc bù đầu đây này.''

Miko chống cằm, đôi mắt híp lại lém lỉnh: ``Hmm\~{} Kuon-senpai thật ra cũng đăng ký rồi, nhưng không nói cho mình biết đó mà, ne\~{}''

``Hay là anh ấy ngại nhỉ?'' Suisei gật gù tán thành.

``Anh ấy chỉ đang giấu sự phấn khích bên trong thôi peko\~{}'' Pekora thêm vào, nửa đùa nửa thật.

Ririka cũng lém lỉnh: ``Có khi anh ấy không đăng ký kịp nên đang làm việc để đỡ cay ấy mà\~{}''

Cả hội xúm vào bàn tán, thì đúng lúc cửa văn phòng mở ra. Subaru bước vào bên trong, giọng oang oang: ``\textbf{Chào buổi trưa cả nhà\~{}!!}''

Fubuki quay phắt lại, gặng hỏi: ``Subaru-chan! Em đăng ký xổ số chưa!?''

Subaru khựng lại, nhíu mày: ``Xổ số gì cơ?''

Không khí trong phòng như đông cứng trong một giây, dường như không một ai tin nổi những gì mà cô vừa thốt ra.

``\textbf{S*itch 2 đó!!}'' Cáo tuyết hét lớn. ``Chị nhắc em mấy lần rồi còn gì! Deadline là đầu giờ chiều nay!''

``Ể-Ể!? Thật á! Em quên mất tiêu!'' Subaru vội vàng móc điện thoại, mở trình duyệt, và rồi\dots\ đứng sững lại.

``Ơ\dots\ nó ghi là 'Hết hạn đăng ký', shuba\dots''

``Hình như là vừa một phút trước rồi đó ạ,'' nàng Thầu khoán lên tiếng, đổ mồ hôi hột với tình cảnh của tiền bối mình.

\dots

\dots

Trong sự im lặng khó tả đó, bất chợt\dots

``\textbf{Trời đất ơi!!!}'' Fubuki ôm đầu ngồi bệt xuống sàn. ``Chị còn nhắn L*NE cho em tận hai lần sáng nay rồi mà!''

Subaru nghe thông tin đó như sét đánh ngang tai. Cô làm rơi điện thoại trên tay, khuỵu người xuống, rồi ngước thẳng lên trần, hét lớn: ``\textbf{Á, shuba!!! Tại sáng nay em lo sửa bài tập voice-over (lồng tiếng) nên quên béng mất\dots!}''

Nỗi đau của người này, trớ trêu thay, lại là niềm vui của người khác. Những người còn lại nổ ra cười ầm. Ririka đập bàn, Miko ôm bụng, Suisei thì lấy điện thoại quay lại biểu cảm chán nản của Subaru. Ngay cả Kaigou, người từ nãy giờ vẫn giữ dáng vẻ nghiêm túc, cũng bật cười không nén được.

Chỉ riêng Kuon vẫn miệt mài trước màn hình. Nhưng ở khóe môi, một nụ cười nhếch mép nhè nhẹ hiện lên, không giấu được chất mỉa mai qua cái chẹp miệng:  ``Đi đăng ký tham dự mà quên cả giờ thì chịu rồi.''

Cậu khẽ lắc đầu, rồi lại quay trở về với tài liệu báo cáo trước mặt, mặc kệ đằng sau là một nàng Vịt đang nằm sõng soài ăn vạ, luôn miệng ``Không công bằng mà, shuba!!'' giữa một đám người vừa thương hại vừa cười hả hê.

\begin{center}
    \uline{Đại học Xây dựng Hà Nội - Phòng 61-H3.}
\end{center}

Buổi chiều hôm đó, khi lên lớp, Yamazu là người đầu tiên bắt chuyện sau khi thấy Kuon và Kaigou vừa đặt cặp xuống bàn: ``\vie{Ủa, nghe nói hai người vừa đăng ký mua S*itch 2 à?}''

Cậu Tổng rướn mày, khẽ ngạc nhiên: ``\vie{Sao ông biết vậy?}''

``\vie{E-Em gái ông đ-đăng tin này lên X mà.}''

Kuon liếc về cô em song sinh, người đang lè lưỡi ra tỏ vẻ vô tội. ``\textit{Có cần thiết phải bắc cái loa phường ra cho mọi người biết không vậy\dots?}''

Cậu bạn Đạo diễn tiếp tục: ``\vie{M-Mà, đến mua máy mà cũng phải rút thăm à? Ghê vãi nhái.}''

Kaigou đáp lại, vẫn còn giữ nguyên sự hào hứng ban sáng: ``\vie{Ừ, ghê đúng không? Năm trăm nghìn máy mà hai triệu lượt đăng ký lận! May là tôi cũng kịp bấm lúc sát nút luôn đó!}''

Tokue từ dãy bên kia, nhìn thấy hai anh em đang nói chuyện sôi nổi, cũng tò mò sán tới: ``\vie{Giớ mới biết là Đại tướng cũng chơi xổ số rinh máy đấy. Nhớ hồi nhập học cũng trúng xổ số chứ? Ba trăm triệu yên\footnote{Xem lại {\flaregothic ÉPISODE 4}, quyển số Không.} gì đó, đúng không nhỉ? Biết đâu ông lại gặp may tiếp?}''

Cậu Tổng gật nhẹ, miệng nửa cười: ``Ờ, đúng là hồi đó có trúng. Nhưng lần này thì\dots\ cũng không hy vọng nhiều. Tờ của tôi có số 4953 lận.''

Nghe đến dãy số của cậu, tất cả mọi người quanh chỗ ngồi cậu đều câm nín.

``\vie{S-Số xui vậy á?}'' Yamazu trố mắt, không tin vào những gì cậu vừa nghe.

``\vie{\dots Ừ, số của tôi là đại hạn cmn luôn.}''

Tamachi ở bên cạnh cậu đẩy gọng kính, nhìn Kuon qua mắt kính mờ phản quang: ``\vie{Dù gì cũng chỉ là con số thôi mà. Ông cũng đừng mê tín quá như vậy chứ.}''

``\vie{Tôi nghĩ câu đó nên dành cho cô em song sinh của tôi thì hợp hơn á,}'' cậu Tổng cười trừ. ``\vie{Tôi thì thế nào chả được. Trúng thì trúng, không trúng thì thôi.}''

Yoisai không ngẩng mặt khỏi laptop, vẫn đang gõ phím: ``\vie{Tôi cũng thấy thế là bình thường. Hồi chơi gacha cũng toàn dính phải mấy số này suốt.}''

Cậu bạn Đeo kính thở ra, ngước về cô em song sinh của cậu: ``\vie{Kaigou thì sao? Số của bà là gì?}''

``8888!'' cô nàng đáp lại.

``\vie{Thế khi của bà có khi lại dính đấy. Số đẹp đến độ tôi ước mua vé số cũng trúng như vậy luôn.}''

Kaigou cười tít mắt: ``\vie{Đúng không!? Biết đâu đấy em lại trúng thật thì sao!}''

Kuon cúi người xuống bàn, chống cằm: ``\vie{Nếu trúng thật thì ít ra cũng có một người trong cái phòng này không cần phải lên phòng thờ cầu xin qua môn nữa.}''

``\vie{Cũng chưa chắc là cậu ta sẽ trúng đâu,}'' Yoisai lên tiếng, sắc mặt nghiêm hơn hẳn. ``\vie{Có vài lần chơi xổ số gacha cũng chẳng trúng tứ quý lần nào.}''

``\vie{Thì thế. Ngẫu nhiên hết mà. Như kiểu mấy cái code anh em mình viết chỉ chạy khi nào nó muốn ấy.}''

Yamazu ngồi cạnh cô bạn Lập trình, lắc đầu bụm miệng cười: ``\vie{M-Mấy ông bà n-nói chuyện máy chơi game mà c-cứ như hội nghị hacker luôn\dots}''

Dẫu như thế nào, một số bạn học đã chúc phúc cho hai anh em nhà Hikawa trúng được xổ số cho chiếc máy này.

\ct

Một tuần trôi qua trong hồi hộp, lo lắng và cả háo hức, giống như họ đang sống lại thời khắc trông ngóng thông báo kết quả dự thi vào đại học.

Đúng bốn giờ chiều giờ Nhật Bản, hòm thư của hàng triệu người trên khắp đất nước bắt đầu ``nổ tung'' khi từng lượt email từ N*nt*ndo được gửi đi, kèm theo một từ định đoạt vận mệnh: 「\ruby{当}{とう}\ruby{選}{せん}」 (Trúng thưởng) hoặc 「\ruby{落}{らく}\ruby{選}{せん}」 (Trượt).

Từ văn phòng \hll\ cho đến các ký túc xá thành viên, đâu đâu cũng vang lên tiếng chuông thông báo, tiếng gào rú phấn khích, và cả tiếng thở dài.

\il

Fubuki là người đầu tiên nổ phát súng ăn mừng trên X với một bài đăng đầy hứng khởi với một từ ngắn gọn:

\txtbx{当選🎉 (tạm dịch: ``Chiến thắng\~{}'')}

Điều bi hài là đến cả người lồng tiếng cho một nhân vật anime nổi tiếng (dường như là người hâm mộ của cô) dỗi tới mức\dots\ ``hít-le'' cô nàng trên X luôn.

Pekora, trong lúc đang trực tiếp trò M*necr*ft, mở mail lên kiểm tra vì khán giả trên bình luận cứ hối thúc nhắc liên tục.

Khi thấy dòng chữ 「当選」, cô nàng suýt làm rớt míc-rô, hét lên: ``\textbf{Uoooooh peko!!!! Trúng rồi peko!!!! S*itch 2 này là của ta!!!!}''

Ngay sau đó, một loạt bình luận chúc mừng cô tràn vào như đê bị vỡ, với đa phần viết ``PEKO WIN!!''

Ngoài ra, một số thành viên khác cũng lần lượt xác nhận đã trúng thưởng:

\begin{itemize}
    \item Okayu bình thản thông báo trong buổi ASMR: ``À, tớ cũng trúng rồi nha. Ngon lành\~{}''
    \item Nene thì đăng một chiếc video nhảy lên vui sướng kèm dòng trạng thái: ``S*itch 2 sẽ là bạn trai mới của tớ!!!''
    \item Luna thì\dots\ lật bàn luôn trên sóng vì vui quá: ``\textbf{Lunaaa dực rùi này\~{}!! Nanora\~{}!!}''
    \item Polka thì mở cả một buổi minishow để ăn mừng chiến thắng của ``tôi và thần may mắn!\footnote{Thực tế, Polka đã trúng chiếc máy ở lượt xổ số thứ hai.}''
\end{itemize}

\il

Thế nhưng, không phải ai cũng may mắn trong trò chơi rút thăm nghiệt ngã này.

Miko, trong một buổi phát sóng trực tiếp, cố gắng giữ vẻ bình tĩnh nhưng cười chát: ``Hề\dots\ Trượt mất rồi ne\~{} Mà thôi, chắc là tại hôm đó mình ấn nút hơi chậm ha\dots''

Suisei không nói gì nhiều, chỉ đăng lên X một câu cụt lủn:

\txtbx{Tui tạch Switch 2 rồi.}

Ririka, ngược lại, thì đầy cảm xúc. Cô đăng một dòng trạng thái với những biểu tượng emoji khóc lóc:

\txtbx{Tui tạch Switch rồi 😭😭😭😭😭}

Những thành viên khác cũng lần lượt báo tin buồn:
\begin{itemize}
    \item AZKi: ``Lỡ hẹn với S*itch 2\dots\ Cũng đành chờ đợt sau vậy\dots''
    \item Ayame thì đăng một bài tweet lên đó, cũng là lần hiếm hoi cô lên mạng: ``Uhm\dots\ rút không trúng\dots\ hơi buồn, nhưng không sao.''
    \item Roboco thì đăng một tấm ảnh được cắt một phần từ thông báo kết quả, với một câu: ``N*nt*ndo\dots\ Ha\dots\ Ha\dots\ Buồn quá😭''
    \item Lamy thì dường như không tin vào mắt mình, chỉ viết một câu: ``Mình không đọc đợc dòng chữ này\dots'' đính kèm một bức ảnh ghi chữ 「落選」to đùng.
\end{itemize}

và còn nhiều thành viên khác nữa. Từ các thành viên cộm cán cho tới những tên lính mới\dots\ tất cả đều chỉ nhận về một cái thông báo đầy phũ phàng.

\il

\begin{center}
    \uline{Đại học Xây dựng Hà Nội - phòng 32-H3.}
\end{center}

Mọi ánh mắt sau đó bắt đầu đổ dồn về hai anh em nhà Hikawa, những người vẫn chưa công bố tình hình.

Hai giờ chiều, chuông báo nghỉ giữa giờ vang lên đúng lúc lớp học vừa chuyển sang phần giải bài tập nhóm. Một số sinh viên tranh thủ ra ngoài mua nước, số còn lại thì lướt điện thoại trong khi chờ giảng viên quay lại, còn một số thì\dots\ đã gục trên bàn từ thuở nào.

Yamazu ngồi ở bên đối diện cậu, bất chợt quay đầu hỏi: ``\vie{Ê, có kết quả xổ số chưa thế?}''

Ngay khi câu hỏi còn chưa kịp lắng xuống, một tiếng thông báo tinh vang lên từ hai chiếc laptop đang đặt cạnh nhau.

``\vie{Ông vừa hỏi xong, có liền luôn,}'' Kuon nhếch môi.

``\vie{Đúng lúc ghê.}'' Kaigou bật nhẹ một tiếng thở, mở hộp thư đến.

Tiếng của cô em vừa dứt, lập tức tất cả các thành viên của nhóm đồ án lập tức sát vào, mắt dán vào hai màn hình máy tính của hai anh em.

Cô em gái song sinh mở mail đầu tiên. ``\vie{Thế nào? Có trúng không?}'' Yoisai hỏi.

Cô đọc lướt qua, rồi bỗng chốc ngả người ra sau ghế, thở ra một hơi dài: ``\vie{Không trúng. Đọc đến `Chúng tôi rất tiếc phải thông báo rằng' là đủ hiểu rồi.}''

Cô bạn Lập trình và cậu Đạo diễn thở dài, như thể họ đã lường trước kết cục này. ``Tiếc ghê.''

``\vie{Mà thôi không sao. Tôi cũng đã đoán trước rồi, không hy vọng nhiều nên không buồn mấy,}'' cô em lên tiếng trấn an mọi người. Dù thế, nét thất vọng vẫn hiện rõ trong đôi mắt của mình.

Giờ chỉ còn mỗi Kuon, ông anh trai thờ ơ với chiếc máy mới. ``Aniki thì sao?''

Lúc này, cậu Tổng vẫn đang kiểm tra hòm thư để xem còn có công việc nào được gửi yêu cầu không, ánh mắt vẫn rất dửng dưng: ``\vie{Chắc là cũng trượt thôi.}''

``\vie{Ô-Ông còn chưa mở t-thì làm sao biết được!}'' cậu bạn Đạo diễn thốt lên, ngạc nhiên với thái độ không mấy quan tâm của cậu. ``\vie{Ít nhất cũng phải xem chứ!}''

Kuon thở dài, liếc sang cậu bạn: ``\vie{Ông trông mong gì từ một cái xổ số có tỷ lệ trúng còn thấp hơn tôi vào Yên Hòa năm tôi đi thi chứ?}''

``\vie{Cứ kiểm tra đi,}'' Yoisai nói. Cô em song sinh của cậu cũng góp giọng: ``\vie{Đúng đó aniki, ít nhất cũng tin vào bản thân chứ!}''

Cậu Tổng thở dài một tiếng, tặc lưỡi: ``\vie{Rồi, kiểm tra thì kiểm tra. Nhưng mà đừng khóc nếu anh trượt nha.}''

Cậu lướt thùng thư tới trên cùng, nhấp vào tin từ N*nt*ndo gửi. Ngay khi tờ thông báo được hiện lên, cả cơ thể cậu khựng lại. ``\vie{V-Vãi\dots}''

Tất cả con mắt đều đổ dồn về phía cậu, khẽ ngạc nhiên với phản ứng bất thường từ cậu. ``\dots Aniki?'' / ``Kuon?'' / ``\vie{Sao thế?}''

Trong bức thư, một khung chữ hiện lên:「当選」, cùng với lời chúc mừng từ công ty.

``\vie{\dots Một phát ăn luôn à?}'' cậu lẩm bẩm. ``\vie{Vãi đạn\dots}''

``\vie{Ơ\dots\ trúng thật à!?}'' Kaigou sững người, đọc lại thông báo một lần nữa.

Khi đã biết chắc anh trai mình đã trúng số, cô nàng mặt bừng sáng, bật dậy khỏi ghế, giọng đầy phấn khích: ``\vie{\textbf{Trúng rồi! Aniki trúng rồi hả!? S*itch 2 về tay nhà mình rồi!!}}''

Tamachi, ngồi ngay cạnh, mỉm cười nhẹ: ``\vie{Chúc mừng ông nhé, Kuon-san.}''

Yoisai quay sang, gật đầu cười: ``\vie{Tỷ lệ trúng thấp thế mà vẫn đậu được. Ông cũng có số thật đó.}''

Tokue chống cằm, khoát tay, cười gian: ``\vie{Thế là định mệnh chọn ông rồi, không thoát được đâu.}''

Shunyo cũng trồi lên: ``\vie{Đại tướng ăn quả này là double win rồi đó! Hết vé số rồi đến chiếc máy. Có gì cho anh em chơi chung nha!}''

Yamazu ngửa mặt ra sau, bật cười: ``\vie{Ghê vãi. Mã số của ông ấy đen thế mà vẫn trúng được. Ô-Ông có vẻ có khiếu chơi lô tô đấy!}''

Kuon thì vẫn chưa nói gì. Cậu nhìn chằm chằm vào màn hình thêm vài giây, rồi khẽ lắc đầu, miệng nhếch thành một nụ cười khó tả, không phải vui sướng, cũng chẳng phải bất mãn, chỉ như thể vừa bị trêu chọc bởi chính số phận.

``Đúng là ghét của nào trời trao của ấy\dots''

Cả nhóm phá lên cười.

N*nt*ndo S*itch 2, tưởng chừng là món quà cho Kaigou, nhưng cuối cùng lại đến tay Kuon. Mà biết đâu, như Tokue nói, đó lại là ``định mệnh''!

\begin{center}
    \uline{Sáng hôm sau\\Văn phòng.}
\end{center}

Không khí trong văn phòng có phần đối lập rõ rệt giữa hai phe: những người chiến thắng và những kẻ thất bại.

Trên những chiếc ghế bành, Miko, Suisei và Ririka đang ngôi tiu nghỉu với những vẻ mặt không khác gì như bị mất sổ gạo. Người thì gục đầu xuống, người thì ôm đầu khóc tu tu, người thì thở dài tới mức không còn năng lượng để nói gì.

Trái ngược lại thảm cảnh đó, do phải đến đầu tháng Sáu, những người trúng xổ số mới có thể mua máy, nhưng niềm vui khi biết chắc chắn mình đã nắm trong tay quyền mua chiếc máy Hai khiến cả Fubuki lẫn Kaigou không giấu nổi vẻ rạng rỡ trên gương mặt.

``Hehe\~{} Chỉ còn đợi vài tuần nữa thôi là được cầm máy thật rồi\~{}'' nàng Cáo tuyết ôm gối, đung đưa chân trên ghế, mắt sáng lấp lánh.

Cô em song sinh thì cầm cốc ca-cao nóng, ngân nga giai điệu chiến thắng, môi khẽ cong lên thành một nụ cười hài lòng: ``Chị thì chắc kèo rằng chiếc máy sẽ nằm trong nhà mình. Thế là đủ rồi.''

``Kuon-senpai thì sao, peko?'' Pekora ngẩng đầu lên khỏi điện thoại, ánh mắt hướng về người đàn ông duy nhất trong căn phòng.''

Kuon vẫn giữ vẻ thản nhiên, như thể chiến thắng của mình là một biến số không đáng bàn tới. Cậu chống cằm, mắt nhìn về phía xa, suy nghĩ về công việc cậu định làm trong ngày. ``Anh thì chẳng quan trọng đâu. Thế nào cũng được,'' cậu buông một câu lạnh nhạt.

``Ồ hố\~{} Thái độ mập mờ này là sao đây\~{}'' Fubuki cất giọng đá xoáy hỏi.

``Không có gì.''

Thật ra, Kuon đã chuẩn bị một lời nói dối tinh vi: giả vờ rằng Kaigou là người trúng thưởng, nhằm tránh mọi nghi ngờ từ các thành viên đã trượt xổ số. Nhưng chưa kịp mở miệng để lên tiếng thì\dots

``À, không phải chị. Là aniki trúng đấy,'' cô em song sinh hồn nhiên chỉ ngón tay về phía cậu.

``Này, Kaigou! Em làm cái quái- Âu sịt.''

Tức thì\dots

Một giây im lặng.

Rồi, như hiệu lệnh vừa được phát đi, ba con người đã ``tạch'' xổ số đồng loạt lao vào như vũ bão, tra hỏi cậu một cách điên loạn như thú dại.

Ririka bùng nổ, hai tay nắm lấy cổ áo, đung đưa cậu qua đi qua lại: ``\textbf{Tại sao anh lại làm thế với em chứ, Kuon-senpai!?}''

``\textbf{Bất công, bất công, bất công!!!}'' Miko nhào tới đòi công bằng, gần như muốn đấm cậu tới nơi.

Suisei ôm đầu, trợn mắt, biểu cảm như đang đứng giữa đêm mưa thất tình: ``\textbf{Trả lại vận may cho em đây!!}''

Nhìn cảnh hỗn loạn giữa ba ``con thú'' đang thèm khát chiếc máy mới cắn xé con mồi, cô em song sinh chỉ cười trừ, miệng mấp máy một lời xin lỗi muộn màng.

Kuon lườm cô em, như thể muốn hỏi: ``Sao lại oang oang mồm ra làm cái gì\dots?''

Nhưng vẫn chưa hết. Thậm chí đến cả hai thành viên đã dành chiến thắng cũng giở giọng châm chọc.

Pekora sau vài giây ngơ ngác cuối cùng cũng hiểu rõ được vấn đề: ``Á à\~{} Quả nhiên Kuon-senpai đúng là đã tham gia xổ số rồi, peko\~{} Vậy mà còn giấu tụi em cơ à\~{}'' giọng cô cao vút, tay chống hông, tai thỏ vểnh lên giận dữ giả vờ.

Fubuki không bỏ lỡ cơ hội châm chọc, miệng cười gian như cáo: ``Ái chà\~{} Kuon-senpai tưởng không thích mà hóa ra vẫn trúng thôi\~{} Cái này người ta gọi là `giả vờ lạnh lùng, ai ngờ bốc trúng máy' đấy\~{}''

Tiếng cười rộn lên khắp phòng làm việc, lấn át luôn cả tiếng điều hòa.

Kuon khẽ thở ra, mặc kệ vẫn đang bị ba cô gái gào khóc kia lắc mạnh như búp bê: ``Tôi đúng là dính nghiệp rồi\dots''

Không ai rõ trong ngày hôm đó Kuon đã giải quyết cho ba cô nàng thua cuộc đó bằng cách nào, chỉ biết rằng, với \hll\, ngay cả một chuyện nhỏ như kết quả xổ số cũng có thể biến thành một trận náo loạn sáng sớm.

\omk

Đầu tháng Sáu, tại một trung tâm thương mại lớn ở Tokyo, khu vực nhận hàng dành riêng cho đợt đầu tiên của chiếc máy Hai đông nghẹt người. Trong số đó, không khó để nhận ra hai bóng dáng quen thuộc bước ra khỏi quầy giao nhận với một túi đựng hộp máy mới toanh: anh em sinh đôi nhà Hikawa.

``\vie{Ồ\dots\ nhìn máy thực tế còn đẹp hơn cả hình mẫu trên mạng nữa,}'' Kaigou ngắm nghía vỏ hộp, đôi mắt sáng như đèn pha. ``\vie{Suzu nói rằng đã rất mong chờ chiếc máy này đó!}''

``\vie{Ừ, công nhận,}'' Kuon gật gù. ``\vie{Chưa kể chiếc hộp này còn nhỏ gọn hơn nữa.}''

Ngay sau đó, một giọng nói hồ hởi vang lên từ phía khác, đến từ một người rất nhẵn mặt.

``A, Kuon-senpai, Kaigou-senpai! Hai người cũng tới nhận máy à?'' Fubuki vẫy tay, cạnh cô là Luna đang líu lo và Okayu với gương mặt điềm đạm, tay ôm túi giấy to.

``Lâu rồi mới thấy hội chiến thắng tụ lại đủ nhỉ,'' nàng Miêu mỉm cười.

``Cúi cùng chúng ta cũng đực tựn tay xở hữu chiếc máy nì rùi đó-nora\~{}'' nàng Công chúa xứ Kẹo ngân nga, cười tít mắt.

Kaigou quay sang chào: ``Lâu lắm mới gặp mấy đứa ở ngoài này! Hôm nay vui thật\~{}''

``Miễn sao mấy đứa thấy hạnh phúc là được,'' Kuon nhún vai, điềm tĩnh đáp.

Tuy nhiên, đối lập vói niềm hân hoan của những người sở hữu được chiếc máy là một đám mây ghen tị từ phía sau ập tới.

Từ phía bên ngoài khu siêu thị, những con người không thể chạm tay vào chiếc máy chơi điện tử cầm tay thế hệ mới nhất đang nhìn nhóm còn lại với những ánh nhìn của một loài thú săn mồi, cùng với những con người đã không kịp đăng ký tham gia xổ số, dù không giận bẳng nhưng cũng không giấu nổi sự ganh ghét.

``Ughhh\dots\ nhìn họ mà tui muốn hóa đá, ne\dots'' Miko chắp tay lên má, mặt nhăn như mèo bị tạt nước.

``Tui cũng muốn Switch 2 màaaa\dots'' Suisei gào trong câm lặng, ánh mắt hướng vào túi của Kuon như thể có tia la-de.

``Mấy người ác quá, ác quá đi mà, shuba\dots'' Subaru lầm bầm, hai tay khoanh lại, cố giữ hình tượng ngầu lòi nhưng khóe môi co giật từng chút một.

``Tôi đâu có tham gia đâu, nhưng tui vẫn ganh tị được nha!'' Korone thản nhiên tuyên bố, gây áp lực tinh thần không kém, đặc biệt là từ cô bạn Mèo tím của mình.

Dẫu vậy, phe chiến thắng không để tâm lắm. Bọn họ đã lên sẵn kế hoạch.

\ct

Khi cả nhóm đang bàn xem nên đi đâu để thử máy, Kaigou hướng đôi mắt mong chờ từ phía ông anh, nhưng Kuon chỉ nhún vai đáp lại: ``\vie{Em trả tiền máy đi.}''

``\vie{Ơ, máy là do anh trúng mà, aniki! Không có lý do gì để em phải trả cả!}'' cô em gái hậm hực phản đối, mặt hơi đỏ lên.

``\vie{Anh trúng, nhưng người dùng là ba đứa bọn em và Miyako. Anh lo phần phụ kiện.}''

``\vie{Hể? Phụ kiện?}''

Cậu rút ra một tờ ghi chú dài hơn biên bản họp, trên đó là những phụ kiện đi kèm. ``\vie{Dock phụ, túi bảo vệ, hai tay cầm Pro Controller (một gốc, một cho GameC*be), miếng dán màn hình, giá đỡ di động, camera. Đã không mua thì thôi, mà đã mua thì phải đủ bộ. Nói trước là giá của chúng cũng đắt không kém gì so với cái máy đâu.}''

``\vie{Vậy chắc là anh em mình\dots\ cam-pu-chia nhỉ?}'' Kaigou gật đầu, thấy rằng người anh nói cũng có lý, dù trong đầu vẫn nghĩ: ``\textit{Thật ra anh chỉ đang trốn tránh việc phải trả tiền mua máy đúng không\dots?}''

Ngay trước khi Kuon quay lại siêu thị để mua phụ kiện, Fubuki bỗng huých nhẹ vào lưng cậu: ``Vậy là bọn em sang nhà anh để test máy được chứ?''

Cậu Tổng nhún vai: ``Ừ, cũng được. Nhưng mà làm ơn đừng có biến nhà anh thành một cái võ đài. Bọn anh vừa mới dọn tối qua xong đấy.''

Tất cả các thành viên còn lại mừng rỡ khi được cậu cho phép. Luna giơ tay reo hò: ``Na\~{} Ta xẽ có tiệc S*itch ở nhà Kuonsa-senpai nanora\~{}''

Và rồi, trong khi phe chiến thắng lặng lẽ rời khỏi hiện trường, mang theo túi hàng nặng trĩu và tâm trạng nhẹ như mây còn Kuon thì quay lại siêu thị mua thêm phụ kiện cho máy mới, thì phe thất bại ở phía sau chỉ còn biết gào rú trong tuyệt vọng.

``\textbf{Tôi ghét cuộc sống này!!!''} Suisei hét giữa đám đông khiến cho những người xung quanh chú ý, thậm chí có người còn bĩu môi cười khi biết được tình hình.

``Đồ phản bội Kuon-senpai!! Em đã từng tin tưởng anh cơ màaa!!'' Ririka đấm không khí.

``Nè Ogayu, bà cũng phải cho tôi chơi đó nha!'' Korone gọi vọng lại từ phía người chiến thắng, và được nàng Mèo đáp lại bằng một ngón cái giơ lên.

Cảnh tượng ấy, nếu được ghi hình, có thể dùng làm quảng cáo cho một bộ phim kinh dị lấy đề tài ``Ganh tị S*itch 2''.

Từ bên trong siêu thị, chứng kiến những con người gầm rú ngoài kia, cậu Tổng chỉ biết đường thở dài: ``\textit{Cuộc đời thật lắm éo le, nhân sâm thì ít, rễ tre thì nhiều.}''

\begin{center}
    \uline{Nhà Hikawa.}
\end{center}

Buổi chiều ngày hôm đó, tầng năm khu A nhà Hikawa rộn ràng hơn thường ngày. Trung tâm của căn phòng là chiếc máy Hai mới tinh của cậu Tổng, kết nối vào TV màn hình lớn, chạy trò chơi vừa ra mắt cùng ngày, M*r*o K*rt World.

Tiếng động cơ, tiếng va chạm xe và âm nhạc đặc trưng của dòng game đua xe vang vọng khắp căn phòng. Kaigou nghiêng người, miệng bặm lại vì hồi hộp khi đang điều khiển nhân vật của cô vượt chướng ngại.

Bên cạnh cô, Okayu vẫn giữ gương mặt điềm nhiên thường lệ nhưng tay bấm rất dứt khoát, trong khi Fubuki thì cười toe: ``Bắt được bà rồi nhé, Okayu\~{}''

Luna ở phía xa lại đang ``cạp'' cần điều khiển, vừa chơi vừa cười khúc khích như trẻ con được quà: ``Vui quá xá-nanora!!''

Sora, trong tư thế đoan trang nhưng ánh mắt thì đầy hưng phấn, lần đầu tiên thử cảm giác điều khiển trên tay cầm mới: ``Hiếm khi mới có dịp chơi cùng mọi người thế này\dots\ Thú vị thật.''

Ở góc phòng, Ryuuhou, Suzuki và Miyako cũng nhập cuộc cùng những tay cầm phụ, chia thành hai đội để thi đấu nhóm. Dù không phải game thủ chuyên nghiệp, cả ba vẫn nhanh chóng bị cuốn vào không khí phấn khích của cuộc đua.

Cô em út của gia đình Hikawa là vui nhất khi cuối cùng cô bé cũng thỏa mãn được ước muốn của mình: ``Game gia đình là nhất! Ryuuhou-onee-sama thấy sao ạ!?''

Cô em thứ thì mỉm cười, nhẹ nhàng đáp lại: ``Quả không bõ công chờ đợi trong hai tháng. Onii-san đúng là số hên thật.''

Miyako thậm chí còn hú lên khi xe của mình lao xuống vực, một điều mà một người kiệm lời như cô rất hiếm khi thể hiện. Rõ ràng là nàng Chó trắng này cũng rất phấn khởi.

Tất cả tạo nên một không gian vừa náo nhiệt, vừa ấm áp,  nơi những con người rất khác biệt cùng tận hưởng niềm vui đơn giản từ một trò chơi. Những khoảnh khắc tông xe nhau, những lời hò reo, thậm chí là những tiếng ca thán thốt lên khi họ thua cuộc\dots\ tất cả đã tạo nên một buổi chiều vui vẻ này, khi họ có chung một niềm vui là được tận hưởng chiếc máy mới.

Nhưng, không ai để ý rằng, trong lúc ấy, Kuon đang ngồi sau cùng, với một chiếc laptop quen thuộc mở trước mặt, và cuộc gọi video giữa cậu và một người đang ở trại quân sự đang bật.

\il

``\vie{Tao nói thật, cái bọn Nhật chúng nó đúng là rảnh vãi lờ,}'' Koheki, cậu bạn chuyên Hóa học cấp ba với Kuon, và cũng là người giúp cậu dựng máy PC đầu tiên từ vài năm trước, lắc đầu ngán ngẩm. ``\vie{Không mua máy ở nước ngoài mà lại đi chơi cái trò xổ số khổ sở đó làm đếch gì?}''

Kuon uống một ngụm cà phê, chẹp miệng một tiếng, rồi thở dài đáp: ``\vie{Tại em gái tao hăng máu thôi. Tao cũng tính là nếu trượt thì sẽ đặt máy ở bên Mẽo hoặc là chờ hàng về Việt Nam rồi. Nhưng mà Kaigou cứ kéo tao vào, và thế là chơi một vé cho vui, ai ngờ lại trúng. Ảo ma vcl luôn.}''

``\vie{Tao còn chẳng biết liệu mày đang số đỏ hay số đen nữa.}''

``\vie{Ừ, tao cũng éo biết đây này,}'' cậu Tổng nhún vai. ``\vie{Trúng xong còn bị đám con gái tẩn hội đồng vì dám giấu chuyện trúng thưởng. May mà họ đánh nhẹ chứ không tao vào Bạch Mai rồi.}''

Cả hai phá lên cười một chút. Sau đó, cậu Hóa học nói nhỏ: ``\vie{À mà\dots\ bọn mày thực ra đâu có sở hữu máy đó đâu nhỉ?}''

Kuon ngả đâu ra sau, thở dài một tiếng như đã biết chác chuyện đó: ``\vie{Thế mới đau chứ. Chỉ có quyền mua và giữ, chứ không phải sở hữu thực sự. Mã máy gắn với tên người trúng. Bọn Nin mà thấy có gì lạ là biến thành cục chặn giấy mười mấy củ ngay, éo nói nhiều.}''

``\vie{Tệ hơn cả hợp đồng cho vay nặng lãi.}''

``\vie{Chuẩn quá.}''

``\vie{Đấy là tao còn chưa nói đến cái giá game cao vô lý nữa đâu. 80 đô cho một game á? Người thường nào kham nổi?}''

``\vie{Nó thực sự éo đáng tiền một chút nào. Cái game kiểu đó bán với giá một nửa là kịch trần rồi. Thế mà vẫn còn có đám Nin con mù quáng mua kia kìa.}''

``\vie{Bọn nó bị tẩy não nặng rồi. Mày quan tâm làm gì cho mệt. Tao có cảm giác rằng mấy đứa em gái mày bị bọn nó làm lu mở lý trí rồi đấy.}''

``\vie{Thế nên tao mới lo đây này\dots}'' Kuon thở dài chán chường. ``\vie{Không sớm thì muộn chúng nó cũng thành nô lệ hết.}''

Một thoáng im lặng, rồi cả hai cùng đồng thanh, như đã tập dượt từ trước: ``\vie{Tóm lại thì\dots\ \emph{PC vẫn là ngon nhất.}}''

``\vie{Mạnh hơn, hình nét hơn,}'' Koheki giơ hai ngón tay làm dấu, ``\vie{không bị khoá vùng, không giới hạn account (tài khoản)}''

``\vie{Lại còn chơi được giả lập,}'' Kuon tiếp lời, mắt thoáng lên vẻ lấp lánh. ``\vie{Từ mấy cái P*2 cùi bắp đến W*i, 3DS, thậm chí cả mấy bản S*itch đời đầu. Chơi mượt, mod tẹt ga, không phải lo lag lủng gì cả.}''

``Tàu hải tặc ra khơi\~{}'' cậu bạn quân nhân huýt sáo. Cậu Tổng cũng vui tính đáp lại: ``\eng{Aye aye, captain!}'' với kiểu chào nhà binh.

Cả hai cậu bạn lại tiếp tục bật cười. Họ thở dài một lúc, rồi chống cằm nói: ``\vie{Tao đang tính nếu mày được nghỉ phép về Hà Nội, tìm thử xem nhà mày còn con card nào mạnh mạnh không. Gắn vào dàn của tao cho đỡ nghẽn cổ chai.}''

``\vie{Ý mày là nâng cấp cho con PC phèn của mày á?}''

``\vie{Ờ. Không có được thì\dots\ giới thiệu cho tao chỗ nào chuyên bán mấy cái card rời mà giá không bị chém lên trời. Hàng second-hand cũng được, miễn là ổn.}''

Koheki gật gù: ``\vie{Rồi. Mày gửi cấu hình dàn cũ cho tao, để tao coi tương thích cái gì. Lúc về tao lục kho test rồi báo giá cho mày.}''

``\vie{OK\~{}}''

Cậu bạn quân nhân ngả người ra ghế, vẻ thư giãn. Kuon thì lại khẽ liếc ra ngoài phòng khách, nơi Kaigou vừa thắng thêm một vòng đua Mario Kart và hét ầm lên vì vui sướng.

``\vie{Chính vì bọn Nin làm ăn bất lương như thế, tao còn đang tính\dots\ tinh chỉnh cái máy S*itch 2 kia một chút nữa.}''

``\vie{Khoan, mày cũng tính hack cái máy luôn à?}''

``\vie{Không đến,}'' cậu Tổng lắc đầu. ``\vie{Chỉ là định cài thêm lớp bảo vệ, kiểu như custom bootloader, hoặc cơ chế rollback (quay về thời điểm an toàn trước), để lỡ có bị bọn Nin dòm ngó thì vẫn giữ được máy. Kiểu chống brick bất công ấy.}''

``\vie{Thế cũng hợp lý,}'' cậu bạn nhướng mày. ``\vie{Tao dám chắc là mày cũng sẽ làm luôn cả phần `tối ưu hóa' hiệu năng, xung nhịp, ép nhiệt, ép điện, tối ưu cache gì đó đúng không?}''

``\vie{Mày hiểu tao rõ thế. Ít nhất cũng phải làm gì với cái máy để xứng đáng chục củ kia chứ,}'' Kuon nhoẻn miệng cười, nửa đùa nửa thật. ``\vie{Tao chỉ giúp mấy em gái nhà tao chơi vui hơn thôi, không để bị bọn kia phá đám.}''

Đến cuối cùng, cả hai cậu bạn đều cùng ngầm hiểu rằng: PC vẫn là thế lực thượng phong trong lĩnh vực trò chơi, còn các máy chơi điện tử khác chỉ dành cho bọn ``trẻ con.''

Kuon gập máy tính, đứng dậy. Trên màn hình TV, nhân vật của em gái song sinh cậu đang nhảy lên bục chiến thắng, còn Fubuki, Luna, Sora và cả Ryuuhou, Suzuki, Miyako thì đang cười vang.

Cậu mỉm cười, rất nhẹ, rồi bước ra nhập cuộc. ``Đến lượt anh chơi chưa nhể?''

``Kuon-senpai cũng muốn nhập hội à?'' Sora cười hiền. ``Vậy bọn em sẽ không chấp anh đâu nha!''

``Zô đi, Kuonsa-senpai!'' Luna vỗ xuống sàn.

``Cẩn thận không là anh ấy lại đẩy bà xuống vực như hồi sinh nhật ba năm trước đó\~{}'' Okayu nói với Fubuki, giọng bông đùa.

``Lần này tôi sẽ phục thù!'' nàng Cáo trắng hạ quyết tâm hết sức.

``Em vẫn còn chấp nhặt chuyện đó à?'' cậu Tổng bật cười.

Nhưng trong lòng cậu, một bản kế hoạch tinh chỉnh phần cứng và nâng cấp dàn máy cũ\dots\ đã bắt đầu khởi động.

\ct

Dù đã có trong tay chiếc máy S*itch 2 từ đợt phát hành đầu tiên, Kuon vẫn không khỏi tò mò dõi theo kết quả những đợt xổ số tiếp theo. Một phần vì em gái song sinh của cậu suốt ngày bàn tán, một phần khác\dots\ là vì trò vui chưa kết thúc.

Dư âm của đợt đầu tiên vẫn chưa nguội thì danh sách trúng thưởng lượt hai đã được công bố. Một sự kiện khiến cộng đồng \hll\ rúng động.

Người chiến thắng là Kanade và\dots\ lại là Polka.

Không ai hiểu bằng cách nào mà nàng Cáo xiếc này có thể trúng tới hai chiếc máy chỉ trong hai lượt. Có người bảo cô lập hai tài khoản, có người thì thở dài: ``Ờ thì hề mà\dots\ đời hề lúc nào cũng bất ngờ.''

Đáng cười nhất là trường hợp của nàng Âm nhạc, mặc dù được quyền mua chiếc máy, cô nàng lại\dots\ không đủ tiền để mua, phải để Nodoka bất đắc dĩ ``giải cứu.''

Trong khi đó, số người trượt thì dài như danh sách cử nhân thất nghiệp. Không thiếu những người đã từng tham gia xổ số đợt trước, và cũng có những cái tên mới: Iroha, Suisei, Noel, Flare, Chihaya\footnote{Ở trong FLOW GLOW. Nhóm này sẽ được giới thiệu ở quyển Hai Mươi Bảy.}, hay Miko, vân vân.

Tất cả đều trắng tay. Mỗi người phản ứng theo cách riêng, nhưng gộp lại thì đều là một bản hợp xướng tuyệt vọng.

``Làm ơn! N*nt*ndo-san, quay lại trúng tôi với!!'' Miko lăn lộn trên sàn.

Suisei thì\dots\ viết liên tục những câu ``Chúng tôi rất tiếc phải thông báo rằng'' tới mức người đọc phát sợ.

``Tôi là Thiên sứ mà còn thua, thế giới này đúng là quá bất công\dots'' Kanata viết status ẩn ý lên X.

\ct

Thêm vài tuần nữa trôi qua. Lượt ba bắt đầu. Thêm một cơ hội, thêm một hy vọng.

Và rồi, những người chiến thắng chiến thắng tiếp theo cũng đã lộ diện: Towa, Nene, và Lamy.

Hai cô nàng thuộc Thế hệ thứ Năm bật khóc, không biết vì vui hay vì sốc. Sau ba lượt mới chạm được vào ánh sáng hy vọng. Nhưng đồng thời, Botan sẽ là người duy nhất bị ``ruồng bỏ'' khi những đồng nghiệp của cô đã sở hữu chiếc máy.

Riêng nàng Ác ma, trúng phát một, bình thản đăng ảnh khoe tờ thông báo trúng thưởng, còn chèn thêm dòng trạng thái: ``Không chơi không trúng, chơi thì trúng luôn\~{}''

Danh sách người trượt lần này dài đến mức đáng được ghi vào sách kỷ lục. Vẫn lại là những gương mặt thân quen, và một số gương mặt mới, có thành viên mới chỉ gia nhập vào công ty được gần một năm.

Nhiều người trong số họ đã trượt đến lần thứ hai, thậm chí thứ ba. Không ít người bắt đầu nghi ngờ vào khái niệm ``xác suất''. Có kẻ đâm ra mê tín, bắt đầu cầu khấn trước khi đăng ký xổ số. Có người lại đổ lỗi cho máy chủ quá khắt khe, hoặc nhà Nin quá tàn nhẫn.

\ct

\pov{Hikawa Kuon}

\begin{center}
    ``\uline{Cuộc thảm sát xổ số S*itch 2 vẫn sẽ còn tiếp tục}''
\end{center}

đó là một sự thật không thể chối cãi, ít nhất là trong khoảng một, hai tháng sắp tới.

Tôi nhấp một ngụm trà, mắt lướt qua chiếc máy Hai sáng bóng trên kệ. Kaigou vẫn đang hí hửng khởi động M*r*o Kart Deluxe trên máy cũ, còn ngoài kia, tiếng rên rỉ thất vọng của những người trượt xổ số vẫn vang lên không ngớt trên các diễn đàn.

Nhìn cảnh đó, tôi không khỏi nhếch mép. Thật nực cười. ``Một cái máy chơi game thôi mà, sao phải biến nó thành một cuộc chiến sinh tồn thế này?'' tôi cá là bạn đang nghĩ vậy.

Nhưng tôi hiểu vì sao mọi người lại phát cuồng với chiếc máy mới này. Phải có lửa thì mới có khói chứ.

Câu trả lời đơn giản thôi: chiêu trò bẩn thỉu của nhà N*nt*ndo. Họ có nhiều chiêu trò làm cho những người dùng như chúng tôi phải khóc dở mếu dở, trong đó giới hạn số lượng và cái xổ số quái quỷ kia là những ví dụ điển hình.

Tôi từng đọc qua về chiến lược của họ: kiểm soát nguồn cung để đẩy giá trị cảm nhận lên cao, khóa tài khoản NSO để ép người chơi vào hệ sinh thái của họ, và rồi còn cái DRM chết tiệt kia, một cơ chế bảo mật khiến máy của bạn có thể biến thành cục gạch nếu bạn dám ``lệch đường''. Chip N*IDIA mới trong Switch 2, nghe đồn là loại cải tiến, có thêm lớp mã hóa để ngăn chặn bẻ khóa. Nhưng tôi cá là nó cũng chẳng bền lâu.

Và đúng như dự đoán, chỉ sau chưa đầy một ngày - đúng, là chưa đến \emph{HAI MƯƠI TƯ GIỜ} - kể từ khi ra mắt, tin tức nổ ra trên các diễn đàn công nghệ. Một nhóm hacker vô danh đã tìm ra cách khai thác ở mức userland - tức là môi trường người dùng, ở bậc thấp nhất. Không phải bẻ khóa hoàn toàn (chưa đâu, cái đó còn lâu) mà chỉ là một bước đầu, cho phép chạy code tùy chỉnh trong không gian người dùng. Với một cách khái thác lỗ hổng ropchain\footnote{Chính tả đúng: ROP chains. Hiểu nôm na là các chuỗi được xây dựng bằng hình thức lập trình hướng kết quả (Return-Oriented Programming). Đây là một kỹ thuật khai thác bảo mật, trong đó hacker sử dụng các đoạn mã ngắn (gadgets) có sẵn trong bộ nhớ của chương trình để thực thi một chuỗi lệnh mong muốn, mà không cần chèn mã mới.}, họ có thể vượt qua một phần hệ thống bảo mật, mở đường cho ``homebrew\footnote{Một phương thức xâm nhập hệ thống bằng trình quản lý gói mã nguồn mở.}'' hoặc custom firmware trong tương lai. Tôi không bất ngờ. Công ty họ N*n có thể khoe chip mới hay DRM xịn, nhưng lịch sử đã chứng minh: không có tường thành nào hacker không leo qua được.

Một số người gọi đó là ``tội phạm công nghệ'', thậm chí còn gọi là ``những tên cướp giữa ban ngày''. Với tôi, đó chỉ là cách cộng đồng lấy lại thứ họ đáng được hưởng: quyền tự do với món đồ họ đã bỏ tiền ra mua.

Tôi không ủng hộ cướp bản quyền, nhưng cái cách mà công ty đó ép người chơi trả 80 đô cho một tựa game, rồi còn khóa vùng và giới hạn tài khoản, khiến tôi ngứa mắt, và tôi hoàn toàn hiểu vì sao họ chọn đi theo con đường ``hắc ám'' đó. Một game như The Legend of Z*l*a mới, dù đẹp đến cỡ nào, cũng không đáng để người thường phải cày cả tháng lương.

Chưa kể, nếu bạn lỡ để máy ``brick'' vì một bản cập nhật lỗi, N*nt*ndo sẽ chẳng buồn cứu bạn. Tôi đã thấy quá nhiều người trên X than thở vì mất cả đống tiền chỉ vì một lỗi phần mềm, thậm chí Ryuuhou cũng đã từng trở thành nạn nhân của lỗi này, và tôi đã mất nguyên một ngày để có thể cứu chiếc máy.

Thế nên, khi tin tức về bản hack lan truyền, tôi chỉ nhún vai. ``\vie{Tự làm tự chịu thôi, Nin à,}'' tôi lẩm bẩm, ``\vie{Các ông cứ làm ăn kiểu này thì đuổi khách người ta đi đấy.}''

Kaigou quay sang, ánh mắt lấp lánh như vừa thắng thêm một vòng đua. ``\vie{Anh lầm bầm gì đấy, aniki? Vẫn cay vì phải mua máy à?}'' cô nàng cười đểu, tay vẫn cầm chặt J*y-C*n.

``Không,'' tôi đáp, giọng tỉnh bơ. ``\vie{Chỉ là đang nghĩ, nếu cái máy này mà bị brick, anh sẽ phải ngồi cả đêm để cứu dữ liệu cho em đấy.}''

Cô nàng phồng má, giả vờ giận. ``\vie{Xì, anh thì lúc nào cũng nghĩ xấu! Cứ tận hưởng đi, aniki! Máy đẹp, game đẹp, thế là đủ rồi!}''

Tôi bật cười, đặt cốc trà xuống. Có lẽ Kaigou đúng. Dù tôi chẳng mê cái máy này, nhưng nhìn lũ em gái và các cô gái \hll\ cười nói rôm rả bên TV, tôi cũng thấy đáng. Dù sao thì, nếu lũ hacker kia tiếp tục phá đảo, có khi tôi cũng sẽ thử cài một custom bootloader cho cái máy này. Chỉ để phòng hờ thôi, biết đâu đấy.

Ở đâu đó ngoài kia, cơn sốt S*itch 2 vẫn đang tiếp diễn, và cả thế giới công nghệ đang chuẩn bị cho một cuộc chiến mới.

Tôi nhún vai, đứng dậy, và bước ra nhập hội với đám đông đang hò hét bên màn hình. ``\vie{Được rồi, tới lượt anh đây. Đừng hòng thắng anh mà dễ, Kaigou.}''

Sau khi lôi được máy Hai về nhà, cô em song sinh của tôi không chỉ nhảy nhót vì vui mà còn mày mò lướt mạng, đọc mấy bài đánh giá từ đám người dùng bên trời Tây. Họ kêu ca đủ thứ: nào là J*y-C*n mới vẫn bị drift (xê dịch), nào là một số game từ máy cũ vẫn chưa tối ưu, chạy giật lag trên máy Hai, thậm chí có người đã phải trả hàng sau vài tiếng mua. Ở Nhật thì ít ai dám nói xấu N*nt*ndo, đất nước sinh ra Nin mà, ai dám ho he? Nhưng Kaigou, dù cuồng máy thật, cũng không hẳn mù quáng, và đó là điều tôi cảm thấy có thể thở phào nhẹ nhõm.

Tôi cũng đã khuyên em ấy nên giữ lại máy Một - chính là máy mà chúng tôi đang chơi đây - vì ít nhất mấy game cũ vẫn chạy mượt mà hơn, và ơn giời là Kaigou đồng ý. Dù sao, mục đích chính của cái máy Hai này là để cả nhà cùng chơi, nhất là Ryuuhou và đặc biệt Suzuki, con bé mê tít mấy game như A-Crossing.

Còn tôi? Tôi vẫn sẽ trung thành với PC và mấy game trên di động, nơi chẳng ai ép tôi phải trả 80 đô hoặc bất kỳ giá nào cao hơn cho một tựa game mà chất lượng bị đặt nhiều nghi vấn. Tôi có thể mua game đường hoàng trên cửa hàng trực tuyến, hoặc không thì ``mượn'' game trên một trang nào khác mà không ai cấm cản gì. Tôi có thể kiên nhẫn chờ hai, ba năm, thậm chí năm hay mười năm để có thể chơi những trò ở trên máy Hai đó ngay trên PC.

Thú thật, tôi cũng hơi bất ngờ khi thấy Kaigou đôi khi lôi PC ra chơi lại một số trò AAA để so sánh, kiểu như muốn kiểm tra xem máy Hai có đáng đồng tiền bát gạo không. Hài thật, cô nàng cuồng Nin là thế, vậy mà vẫn còn chút tỉnh táo để phân tích. Có lẽ cô em tôi chưa bị N*nt*nd* dắt mũi hoàn toàn và tôi nên thấy mừng vì điều đó.

Nếu mọi thứ vẫn diễn ra như vậy, nhà Nin sẽ trở thành một tiền lệ cực xấu để các ông lớn trong làng trò chơi điện tử noi theo, vốn cũng là những chất xúc tác dần phá hủy ngành công nghiệp trò chơi điện tử, và cuối cùng\dots

\dots sẽ chỉ dẫn đến một năm 1983 khác ngay tại thế kỷ XXI này, không sớm thì muộn.

\begin{center}{\abv\Large Bạn có biết?}\end{center}

Tính tới tháng Sáu năm 2025, sau bốn lượt xổ số, dưới đây là danh sách những người thắng và thua:

\begin{center}
    \begin{tabular}{|c|>{\raggedright\arraybackslash}p{0.4\textwidth}|>{\raggedright\arraybackslash}p{0.4\textwidth}|}
        \hline
               & Chiến thắng                       & Thua cuộc                                                                                                                                                              \\
        \hline
        Lượt 1 & Pekora, Fubuki, Okayu, Luna, Sora & Miko, Suisei, AZKi, Roboco, Ayame, Mio, Noel, Flare, Kanata, Towa, Botan, Watame, Lamy, Lui, Koyori, Ririka, Chihaya, [Korone, Subaru, La+ và Kanade] (không tham gia)
        \\
        \hline
        Lượt 2 & Kanade, Polka                     & Iroha, Suisei, Noel, Koyori, Lamy, AZKi, Nene, Ayame, Ririka, Kanata, Hajime, Lui, Watame, Roboco, Flare, Chihaya, Miko                                                \\
        \hline
        Lượt 3 & Towa, Nene, Lamy                  & Koyori, Noel, Miko, Suisei, Lui, Ririka, Roboco, Flare, Mio, Ayame, AZKi, Iroha, Watame, Botan, Chihaya, Kanata, Su, Hajime, Subaru                                    \\
        \hline
        Lượt 4 & Ayame, Iroha                      & Suisei, Ririka, Flare, AZKi, Su, Kanata, Subaru, Miko, Botan, Lui, Koyori, Hajime, Noel, Chihaya, Mio, Watame                                                          \\
        \hline
    \end{tabular}
\end{center}

\end{document}